\documentclass[11pt]{article}

\usepackage[utf8]{inputenc}
\usepackage[english]{babel}
\usepackage{enumerate}
\usepackage{latexsym}
\usepackage{amsthm,amsmath}
\usepackage{amssymb}
\usepackage{aliascnt} %para que funcione bien \Cref
\usepackage{multicol}
\usepackage{tikz}
\usepackage{float} % for [H] option in figures


\frenchspacing
\usepackage[bookmarks]{hyperref} 
\hypersetup{
	pdftitle={Classification of prime graphs with 2-switch-degree at most 4},
	pdfauthor={Victor N. Schvollner},
	colorlinks=true, linkcolor=blue, citecolor=red,
	filecolor=cyan, urlcolor=magenta
}

\usetikzlibrary{babel,decorations.pathreplacing,calc,positioning,shapes.geometric,arrows.meta}

\usepackage{cleveref}

% Color definitions remain unchanged
\definecolor{10}{RGB}{115,59,171}
\definecolor{8}{RGB}{212,122,240}
\definecolor{7}{RGB}{99,212,119}
\definecolor{6}{RGB}{183,240,164}
\definecolor{D}{RGB}{255,162,79}
\definecolor{E}{RGB}{255,84,0}
\definecolor{F}{RGB}{158,248,255}
\definecolor{G}{RGB}{128,135,255}
\definecolor{I}{RGB}{187,255,0}
\definecolor{A}{cmyk}{.9,.05,.4,0}
\definecolor{B}{RGB}{150,30,150}
\definecolor{C}{RGB}{186,155,189}
\definecolor{9}{RGB}{0,180,60}
\definecolor{0}{RGB}{30,123,191}
\definecolor{1}{RGB}{255,113,102}
\definecolor{2}{RGB}{41,199,92}
\definecolor{3}{RGB}{242,207,16}
\definecolor{5}{RGB}{255,15,154}
\definecolor{4}{rgb}{.8,0,.8}

\definecolor{Red}{rgb}{1,0.4,0.4}
\definecolor{Green}{rgb}{.1,.5,.1}
\definecolor{Blue}{rgb}{.1,.1,.5}
\definecolor{blue}{RGB}{0,0,255}
\definecolor{Yellow}{rgb}{.8,.4,0}
\definecolor{X}{rgb}{.8,.4,0}
\definecolor{H}{rgb}{0,0,1}
\definecolor{light}{rgb}{.67,.84,.90}
\definecolor{Cyan}{rgb}{0,1,1}
\definecolor{Purple}{rgb}{.5,0,.5}
\definecolor{Purple2}{rgb}{.5,.2,.5}
\definecolor{white}{rgb}{1.0,1.0,1.0}
\definecolor{Purple2}{rgb}{.8,.4,0}
\definecolor{Amarillo}{RGB}{225,191,73}
\definecolor{Celeste}{RGB}{117,170,219}
\definecolor{Castano}{RGB}{232,53,17}
\definecolor{Black}{RGB}{0,0,0}
\definecolor{White}{RGB}{255,255,255}
\definecolor{gris}{rgb}{.5,.5,.5}

\DeclareMathOperator{\act}{act}


\newtheorem{theorem}{Theorem}[section]

\newaliascnt{corollary}{theorem}%para que funcione bien \Cref
\newtheorem{corollary}[corollary]{Corollary}
\aliascntresetthe{corollary}%para que funcione bien \Cref

\newaliascnt{lemma}{theorem}%para que funcione bien \Cref
\newtheorem{lemma}[lemma]{Lemma}
\aliascntresetthe{lemma}%para que funcione bien \Cref

\newaliascnt{proposition}{theorem}%para que funcione bien \Cref
\newtheorem{proposition}[proposition]{Proposition}
\aliascntresetthe{proposition}%para que funcione bien \Cref

\newcommand{\ap}[1]{\textcolor{red}{\sf{#1}}}

% TikZ Layers
\pgfdeclarelayer{background2}
\pgfdeclarelayer{background}
\pgfdeclarelayer{foreground}
\pgfsetlayers{background2,background,main,foreground}


\title{Complete split and non-split prime graph classification at 2-switch-degree four}
\author{Victor N. Schvöllner}
\date{arxiv:2609.40274v1; CC BY4.0}
\begin{document}
\maketitle
\noindent\textit{Editorial note.} One complete prime-degree-four classification, combining original split and non-split theorems. Complete original mathematical paper138--2741, all auxiliary structural discussion, every active native inline graph figure and embedded bibliography retained verbatim. Original split theorem proof refers to preceding discussion; that complete discussion is included, and no new proof environment or argument is invented.\par\medskip
\noindent\textit{Editorial note.} Lower-degree classifications, realization-graph consequences, global transitivity/regularity dichotomy and counts of derived realization graphs remain context, not additional counts.\par\medskip
\noindent\textit{Editorial note.} Author retains responsibility for mathematics. No proof repair, new assumptions, independent mathematical review, graph enumeration/counterexample search or certification.\par\medskip
\noindent\textit{Allowed standard prerequisites (editorial).}\begin{itemize}
\item Degree-sequence realization graphs, 2-switches, active vertices, split graphs and Tyshkevich composition. Original active/primality/indecomposability/switch-degree definitions and graph-isomorphism/complement conventions retained. Cited realization-graph decomposition results remain original background.. Complete author local-counting and structural case discussion retained, including all inline figures; no graph enumeration, independent counterexample search or author proof certification.
\end{itemize}\medskip
\section{Introduction}

 

The systematic study of the graphs sharing a common degree sequence has
long been a central theme in graph theory. A fundamental tool in this
context is the \emph{2-switch}: given two vertex-disjoint edges $ab$ and
$xy$ of a graph $G$ with $ax,by\notin E(G)$, the 2-switch on
$\{ab,xy\}$ replaces these two edges by $ax$ and $by$
(see \cite{chartrand2010graphs}). 
Since it preserves degrees, it moves between the
realizations of a fixed degree sequence $d$, which form the vertices of
the \emph{realization graph} $\mathcal{G}(d)$: the labeled graphs on a
fixed vertex set with degree sequence $d$, two being adjacent whenever
one is obtained from the other by a single 2-switch
(see \cite{arikati1999realization}). Classically, any two realizations of
$d$ are related by a finite sequence of 2-switches (see page 24 of
\cite{chartrand2010graphs}); equivalently, $\mathcal{G}(d)$ is connected.

Understanding the structure of $\mathcal{G}(d)$ is a rich area of
research: Arikati and Peled proved in \cite{arikati1999realization} that $\mathcal{G}(d)$ is Hamiltonian
whenever $d$ has majorization gap~1,
Taylor studied in \cite{taylor2006contrained} constrained versions of the switching operation, and more recent work has focused on the
connectivity of the subgraphs of $\mathcal{G}(d)$ induced by particular
graph families (see  \cite{schvollner_pseudoforests}).
A more granular, local description of $\mathcal{G}(d)$ is provided by
the \emph{2-switch-degree} (or simply the \emph{degree}) of a graph $G$,
denoted $\deg(G)$, which is the number of distinct 2-switches that can
be performed on $G$. Since two distinct 2-switches on $G$ produce two
distinct graphs, $\deg(G)$ is precisely the degree of $G$ as a vertex of
$\mathcal{G}(d)$. This parameter was introduced
in \cite{pastine20252}, where it is expressed in terms of the induced
subgraphs of $G$ of order~4.

The study of 2-switch transformations is closely related to the
Tyshkevich canonical decomposition \cite{tyshkevich2000decomposition}, which expresses every graph as a
composition $G=G_n\circ\cdots\circ G_1$ of indecomposable graphs, in a
way that is unique up to isomorphism when all the factors have at least one vertex. A key feature of this decomposition
is that all its factors, with the possible exception of $G_1$, are split
graphs. A graph $S$ is \emph{split} if its vertex set can be
partitioned into a clique $K$ and an independent set $I$; when such a
partition is fixed we write
$(S,K,I)$ and we say that $(K,I)$ is a \emph{bipartition} for $S$. Split graphs arise naturally in several areas of graph theory
and have been widely studied
(see \cite{hammer1977split, splitnordhausgaddum, whitman2020split,jaume2025nullspace}). For split graphs the 2-switch-degree admits a
particularly simple description: it equals the number of induced
$P_4$'s.

A vertex of $G$ is \emph{active} if it takes part in some 2-switch on $G$. A graph whose vertices are all active is itself called \emph{active}, and an indecomposable active graph is called \emph{prime}
(see \cite{pastine20252}). Since every factor in the Tyshkevich decomposition
of an active graph is prime, and since all these factors but at most one
are split, prime split graphs are the basic building blocks of the whole
theory. Their structure is further clarified by the \emph{factor graph}
$\Phi=\Phi(S,K,I)$, a loopless multigraph on the vertex set $I$ that records the
2-switches available in $S$. This notion was introduced
in \cite{pastine2025simple}, inspired by the $A_4$-structure of Barrus
and West \cite{barrus.west.A4}; there it is shown that an active split
graph $S$ is prime if and only if $\Phi$ is connected.

Building on these foundations, we classify the graphs whose
2-switch-degree is small, that is, those that are least flexible under
2-switch transformations. Our starting point is a complete classification
of the active graphs of degree at most~3
(\Cref{clasificacion.activo_deg=2,clasificacion.activo_deg=3}); the prime
ones among them are all split except $C_4$ and $2K_2$. This already
determines the global shape of a realization graph from a single one of
its local degrees: if $X$ realizes $d$ and $\deg(X)=k\le3$, then
$\mathcal{G}(d)$ is vertex-transitive and $k$-regular if and only if
$T_{221}$ and $\overline{T_{221}}$ are not induced subgraphs of $X$
(\Cref{G(d).transitivo}), the only exception up to complementation. The
hypothesis $k\le3$ cannot be dropped: for every $k\ge 4$ there is a prime
graph $G_k$ of degree $k$ carrying a 2-switch $\tau$ with
$\deg(\tau(G_k))=2k-2$ (\Cref{familia.Gk}), so from degree~4 on a single
2-switch may almost double the degree and regularity fails. Degree~4 is
thus the first value at which the picture is no longer rigid, and it is
the one we settle here: there are exactly 18 prime graphs of degree~4
(\Cref{clasificacion.primos.deg=4}), 12 split
(\Cref{clasificacion.split.primos.deg=4}) and 6 non-split
(\Cref{clasificacion.no.split.primos.deg=4}). Finally, we determine the
realization graphs of all 29 prime graphs of degree at most~4
(\Cref{realization.graphs.deg<5}): up to isomorphism there are only 10, so
a short list accounts for every realization graph arising in this range.
Some of our tools are stated beyond degree~4; for instance,
\Cref{C4.distinguidores} applies to every graph of degree at most~5.

The rest of the article is organized as follows. \Cref{sec_tools}
gathers the tools used throughout the paper. In \Cref{sec_deg<3} we
classify the active graphs of degree at most~2, and in \Cref{sec_deg=3}
those of degree~3. \Cref{sec_deg=4,sec_deg=4.no.split} are devoted to the
prime graphs of degree~4, split and non-split respectively. \Cref{sec_realization} studies the regularity of
realization graphs having a vertex of degree $k\le4$ and classifies them.

%%%%%%%%%%%%%%%%%%%%%%%%%%%%%%%%%%%%%%%%%%%%%%%%%%%%%%%%%%%


\section{Tools}\label{sec_tools}	

Let $G$ be a graph. We write $V(G)$ and $E(G)$ for its vertex set and
its edge set, and $|G|=|V(G)|$ for its order. The complement of $G$ is
denoted by $\overline{G}$. We write $H\prec G$ to mean that $H$ is an
induced subgraph of $G$, and $G[W]$ (or $G[v_1,\dots,v_k]$) for the
subgraph of $G$ induced by $W\subseteq V(G)$. The notation $G\approx H$
means that $G$ and $H$ are isomorphic. For $n\in\mathbb{N}=\{1,2,\ldots\}$ we put
$[n]=\{1,\dots,n\}$, and $\dot\cup$ denotes disjoint union.
The \emph{neighborhood} of $v\in V(G)$ is
$N_G(v)=\{x\in V(G):vx\in E(G)\}$, and $\deg_G(v)=|N_G(v)|$ is the
\emph{degree} of $v$ in $G$. When $G$ is clear from the context we
write $N_v$ and $d_v$. A path is denoted by the juxtaposition of its
vertices, so $uxyv$ is the $P_4$ with edges $ux,xy,yv$; for a cycle, the first vertex is repeated at the end. We write
$\omega(G)$ and $\alpha(G)$ for the clique number and the independence
number of $G$. 
If $V(G)=[n]$, the \emph{degree sequence} of $G$ is
$d=d(G)=(d_1,\dots,d_n)$. If $d=(d_v)_{v=1}^n$ is the degree
sequence of $G$, then $\overline{d}=(n-1-d_v)_{v=1}^n$ is the degree
sequence of $\overline{G}$. For graph-theoretical concepts not defined here we refer the reader to
\cite{chartrand2010graphs}.

For graphs $G$ and $H$, we denote by $\mathcal{Q}_G(H)$ the set of induced
subgraphs of $G$ isomorphic to $H$, and we put $\mathcal{Q}_G=\mathcal{Q}_G(P_4)\,\dot\cup\,\mathcal{Q}_G(C_4)\,\dot\cup\,\mathcal{Q}_G(2K_2)$.
If $\tau$ is a 2-switch on $G$, we write $\tau(G)$ for the resulting
graph. The four vertices on which $\tau$ acts induce a member of
$\mathcal{Q}_G$, and the next result counts the 2-switches on $G$ accordingly.

\begin{theorem}[\cite{pastine20252}]
	\label{degreeofG}
	For every graph $G$,
	\begin{equation}
		\label{eq18}
		\deg(G)=2|\mathcal{Q}_G (2K_2 )|+2|\mathcal{Q}_G (C_4 )|+|\mathcal{Q}_G (P_4 )|.
	\end{equation}
	In particular, if $G$ is a split graph, then $\deg(G)=|\mathcal{Q}_G (P_4 )|$.
\end{theorem}

By \eqref{eq18}, $\deg(P_4)=1$ and $\deg(C_4)=\deg(2K_2)=2$, so
\eqref{eq18} can be rewritten as
\begin{equation}
	\label{deg.como.suma}
	\deg(G)=\sum_{H\in \mathcal{Q}_G}\deg(H).
\end{equation}



%\Cref{degreeofG} allows us to easily characterize graphs $G$ of degree $\leq 2$ based on the members of $\mathcal{Q}_G^*$.
%%\begin{corollary}
%If $G$ is a graph with $|G|>0$, then the following statements are obviously equivalent to say that $\deg(G)=0$:
%	\begin{enumerate}
%		\item $G$ has no induced subgraphs isomorphic to $2K_2$, $C_4$, or $P_4$ (it follows from \Cref{degreeofG});
%		\item $G$ is threshold (threshold graphs are precisely those which contain no induced subgraphs isomorphic to $2K_2$, $C_4$ or $P_4$; see \cite{mahadev1995threshold});
%		%\item $G$ is inactive (see \cite{pastine20252});
%		\item $\act(G)=\varnothing$.
%	\end{enumerate}
%%\end{corollary}
%
%%\begin{corollary}
%	A graph $G$ has degree 1 if and only if $G$ has exactly one induced subgraph isomorphic to $P_4$ and no induced subgraphs isomorphic to $2K_2$ or $C_4$.
%%\end{corollary}
%
%%\begin{corollary}
%	A graph $G$ has degree 2 if and only if exactly one of the following conditions holds:
%	\begin{enumerate}
%		\item $G$ has exactly two induced subgraphs isomorphic to $P_4$ and no induced subgraphs isomorphic to $2K_2$ or $C_4$.
%		\item $G$ has exactly one induced subgraph isomorphic to $C_4$ and no induced subgraphs isomorphic to $2K_2$ or $P_4$.
%		\item $G$ has exactly one induced subgraph isomorphic to $2K_2$ and no induced subgraphs isomorphic to $C_4$ or $P_4$.
%	\end{enumerate}
%\end{corollary}

\begin{theorem}
	\label{|G^*|}
	If $G$ is an active graph, then $4 \leq |G| \leq 4\deg(G)$.
\end{theorem}

\begin{proof}
	Since $G$ is active, it is clear that $\mathcal{Q}_G \neq \varnothing$, as $\deg(G) \geq 1$. Thus, $|G| \geq 4$. On the other hand:
	\begin{equation*}
		|G| = \left|\bigcup_{H \in \mathcal{Q}_G} V(H)\right| \leq \sum_{H \in \mathcal{Q}_G} 4 \leq 4 \sum_{H \in \mathcal{Q}_G} \deg(H) = 4\deg(G),
	\end{equation*}
	the last equality by \eqref{deg.como.suma}.
\end{proof}



%\begin{corollary}
%	If $d$ is a degree sequence, then
%	\begin{equation*}
%		|\act(d)| \leq 4\min\{\deg(G) : G \in V(\mathcal{G}(d))\}.
%	\end{equation*}
%\end{corollary}

 
%As shown in \cite{pastine20252}, $\act(G)$ depends only on the degree sequence $d$ of $G$, so we may write $\act(d)$.


\subsection{Split graphs and factor graphs}

If $(S,K,I)$ is a split graph and $G$ is a graph disjoint from $S$, the
\emph{Tyshkevich composition} $(S,K,I)\circ G$ is the graph with vertex
set $V(S)\cup V(G)$ and edge set
$E(S)\cup E(G)\cup\{xy:x\in K,\ y\in V(G)\}$
(see \cite{tyshkevich2000decomposition}); we denote by $S^n$
the composition of $n$ disjoint copies of a prime split graph $S$. A graph $G$ is
\emph{decomposable} if $G=S\circ H$ for some graphs $S$ and $H$ with $|S|,|H|\geq 1$, and \emph{indecomposable} otherwise.

The \emph{factor graph} of a split graph $(S,K,I)$ is the loopless
multigraph $\Phi=\Phi(S,K,I)$ with $V(\Phi)=I$ having one edge between $u$
and $v$ for each induced $P_4$ of $S$ containing both $u$ and $v$
\cite{pastine2025simple}. The multiplicity of the edge $uv$ is denoted by
$\sigma_{S}(uv)$, or simply $\sigma_{uv}$. As shown in
\cite{pastine2025simple}, the multiset $E(\Phi)$ does not depend on the bipartition, so the notation $\Phi(S)$ is unambiguous when $S$ is active, and $\deg(S)=|\mathcal{Q}_S(P_4)|=|E(\Phi(S))|$. 
%Every induced $P_4$ of $S$ has its two ends in $I$ and
%its two inner vertices in $K$ \cite{pastine2025simple}; hence
%, by \Cref{degreeofG}.

\begin{theorem}[\cite{pastine2025simple}]
	\label{S.primo.iff.Phi(S).conexo}
	If $S$ is an active split graph, then $S$ is prime if and only if $\Phi(S)$ is connected. 
%	\begin{enumerate}
%		\item 
%		
%		\item If $S=S_n\circ\cdots\circ S_1$, where each $S_k$ is prime, then $\Phi(S)=\dot{\bigcup}_{k=1}^n\Phi(S_k)$.
%		%\item If $G$ is an active graph, then each factor in the Tyshkevich decomposition of $G$ is prime.
%	\end{enumerate}
\end{theorem}

\begin{proposition}
	\label{|I|_|K|<=2deg(S)}
	If $(S,K,I)$ is a prime split graph, then 
	\[2 \leq |I|, |K| \leq \deg(S)+1.\]
%	\begin{enumerate}
%		\item If $S$ is active, then $2 \leq |I|, |K| \leq 2\deg(S)$.
%		
%		\item If $S$ is prime, then $2 \leq |I|, |K| \leq \deg(S)+1$.
%	\end{enumerate}
\end{proposition}

\begin{proof}
	Every member of $\mathcal{Q}_S(P_4)$ meets $I$ in exactly two vertices, and likewise $K$ (see \cite{pastine2025simple}). If $|I| \leq 1$ or $|K| \leq 1$, then $\mathcal{Q}_S(P_4) = \varnothing$, so $\deg(S)=0$ by \Cref{degreeofG} and no vertex of $S$ is active. This gives $2 \leq |I|,|K|$.
%	\begin{enumerate}[(1).]
%		\item Since $S$ is active, every vertex of $S$ lies in some member of $\mathcal{Q}_S(P_4)$. Therefore, 
%		\[ |I| = \left|\bigcup_{H \in \mathcal{Q}_S(P_4)} (V(H) \cap I)\right| \leq\] \[ \sum_{H \in \mathcal{Q}_S(P_4)} |V(H) \cap I| = 2|\mathcal{Q}_S(P_4)| = 2\deg(S), \]
%		the last equality by \Cref{degreeofG}. Replacing $I$ by $K$ gives $|K| \leq 2\deg(S)$.
		%\item 
		
		By \Cref{S.primo.iff.Phi(S).conexo}, $\Phi(S)$ is connected. Since $V(\Phi(S))=I$, $|E(\Phi(S))|=\deg(S)$, and a connected multigraph of size $m$ has order at most $m+1$, we obtain $|I| \leq \deg(S)+1$. Now $(\overline{S},I,K)$ is a prime split graph with clique $I$ and independent set $K$ (see page 14 of \cite{tyshkevich2000decomposition} and page 5 of \cite{pastine20252}). Thus, the previous argument applied to it yields $|K| \leq \deg(\overline{S})+1 = \deg(S)+1$ (see \cite{pastine20252} for the last equality). 
	%\end{enumerate}
\end{proof}
	%If $S \approx P_4^{\deg(S)}$, it is evident that $|I|=|K|=2\deg(S)$. For the converse, we proceed by induction on $n=\deg(S)$. Clearly, $|I|=|K|=2$ implies $S \approx P_4$. If $n \geq 2$, assume the following: if $(R,K',I')$ is an active split graph such that $|I'|=|K'|=2(n-1)$, then $R \approx P_4^{n-1}$. If $H \in \mathcal{Q}_S(P_4)$, consider the split graph $(R,K',I')=S-V(H) \prec S$. Since $|I|=|K|=2n$, we have $V(H_i) \cap V(H_j) = \varnothing$ for $i \neq j$ and $\{H_i,H_j\} \subseteq \mathcal{Q}_S(P_4)$. Thus, $|I'|=|K'|=2(n-1)$, and by the inductive hypothesis, $R \approx P_4^{n-1}$. We need to prove that $E(S)=E(R \circ H)$...
%\end{proof}

\begin{lemma}
	\label{5>degG=|{P_4}|.implica.G.split}
	If $G$ is a graph such that $\deg(G)=|\mathcal{Q}_G(P_4)| \leq 4$, then $G$ is split.
\end{lemma}

\begin{proof}
	It is a well-known fact that a graph is split if and only if it contains no induced subgraphs isomorphic to $C_4$, $2K_2$, or $C_5$ (see \cite{hammer1977split}). The hypotheses imply that $G$ contains no induced subgraphs isomorphic to $C_4$ or $2K_2$. If $C_5 \prec G$, then $4 \geq \deg(G) \geq \deg(C_5)=5$.
	%, by  \Cref{deg.respects.ind.inc}. 
	Thus, $G$ cannot contain induced subgraphs isomorphic to $C_5$ either.
\end{proof}

%\begin{theorem}[\cite{pastine20252}]
%	\label{SoG.act.iff.S_G.act}
%	$S\circ G$ is active if and only if $S$ and $G$ are active.
%\end{theorem}

%\begin{corollary}[\cite{pastine20252}]
%	\label{degG=deg(G.complemento)}
%	For any graph $G$,
%	\begin{equation*}
%		\deg(G)=\deg(\overline{G}).
%	\end{equation*}
%\end{corollary}

A split graph $S$ is indecomposable if and only if $\overline{S}$ is (see
\cite{tyshkevich2000decomposition}, page 14); a vertex is active in $S$ if
and only if it is active in $\overline{S}$ (see page 5 of \cite{pastine20252}),
and $\deg(\overline{S})=\deg(S)$ (see Corollary 6.2 of \cite{pastine20252});
and if $(K,I)$ is the bipartition of $S$, then $(I,K)$ is that of
$\overline{S}$, so $\omega(\overline{S})=\alpha(S)$ and
$\alpha(\overline{S})=\omega(S)$. Hence, to determine the prime split graphs
of degree $k$, it suffices to find those with $\alpha(S)\leq\omega(S)$, i.e.,
$|I|\leq|K|$, and add the complements of those with $\alpha(S)<\omega(S)$.
Those with $\alpha(S)=\omega(S)$ need not be self-complementary: for
instance, $R_{211}$ satisfies $\alpha=\omega=3$ and
$\overline{R_{211}}\not\approx R_{211}$ (see \Cref{split.primos.deg=4.fig.3}).


The following proposition lists some basic properties of $\sigma_{uv}$.

\begin{proposition}[\cite{pastine2025simple}]
	\label{prop.basicas.sigma_uv}
	Let $S$ be a split graph.
	\begin{enumerate}
		\item $\sigma_{uv}=(d_u-\eta_{uv})(d_v-\eta_{uv})$, where $\eta_{uv}=|N_u\cap N_v|$;
		
		\item $\sigma_{uv}=0$ and $d_v \leq d_u$ if and only if $N_v \subseteq N_u$; 
		
		\item $\sigma_{uv}=0$ and $d_u =d_v$ if and only if $N_u = N_v$;
		
		%\item if $N_u = N_v $, then $N_{\Phi}(u)=N_{\Phi}(v)$;
		
		\item if $\sigma_{uv}=1$, then $d_u =d_v$ and $|N_u -N_v |=|N_v -N_u |=1$;
		
		\item if $d_u \geq d_v$ and $\sigma_{uv}=p$ is prime, then $|N_u -N_v |=p, |N_v -N_u |=1$ and $d_u-d_v=p-1$;
		%\item $d_u =d_v$ if and only if $|N_u-N_v|=|N_v-N_u|$;
		
		%\item if $d_u =d_v$, then $\sigma_{uv}$ is a perfect square;
		
		%\item if $S$ has no isolated vertices, then $\sigma_{uv}=d_u d_v$ if and only if $N_u \cap N_v =\varnothing$.
	\end{enumerate}
\end{proposition}

From now on, keep in mind that isolated and universal vertices of a graph $G$ are not active vertices of $G$ (see \cite{pastine20252}).

\begin{lemma}
	\label{S_activo_|I|=2}
	Let $(S,K,I)$ be an active split graph with $I=\{a,b\}$. Then: 
	\begin{enumerate}
		\item $|K|=d_a+d_b$ and $\sigma_{ab}=d_ad_b$;
		
		\item $S$ is isomorphic to the prime split graph with bipartition $([|K|],I)$ such that $N_a=[d_a]$ and $N_b=\{d_a+1,\ldots,|K|\}$, which we call $D_{d_b,d_a}$.
	\end{enumerate}
\end{lemma}

\begin{proof}
	\begin{enumerate}
		\item Since $S$ has no universal vertices, we have $\eta_{ab}=0$ and so $|K|=d_a+d_b$. Moreover, $\sigma_{ab}=d_ad_b$ by \Cref{prop.basicas.sigma_uv}.
		
		\item Straightforward. \qedhere
	\end{enumerate}
\end{proof}

\begin{lemma}
	\label{S_activo_|I|=3}
	Let $(S,K,I)$ be an active split graph with $I=\{a,b,c\}$. For $u\in I$, let $\pi_u$ be the number of vertices of $K$ whose only neighbor in $I$ is $u$, and let $\kappa_u$ be the number of vertices of $K$ whose neighbors in $I$ are exactly the two vertices of $I-\{u\}$. Then:
	\begin{enumerate}
		\item $|K|=\sum_{u\in I}(\pi_u+\kappa_u)$ and $d_u=\pi_u+\sum_{v\in I-\{u\}}\kappa_v$ for all $u\in I$;
		\item $\sigma_{uv}=(\pi_u+\kappa_v)(\pi_v+\kappa_u)$ for all distinct $u,v\in I$;
		\item $S$ is determined, up to isomorphism, by $(\pi_a,\pi_b,\pi_c,\kappa_a,\kappa_b,\kappa_c)$;
		\item if $\sigma_{ac}=0$, then, up to exchanging $a$ and $c$, we have $\pi_a=\kappa_c=0$, $\sigma_{ab}=\kappa_b(\pi_b+\kappa_a)$, $\sigma_{bc}=\pi_b(\pi_c+\kappa_b)$ and $|K|=\pi_b+\pi_c+\kappa_a+\kappa_b$.
	\end{enumerate}
\end{lemma}

\begin{proof}
	Since $S$ is active, every $x\in K$ is an inner vertex of an induced $P_4$ of $S$, whose ends lie in $I$ (see \cite{pastine2025simple}), so $N_x\cap I\neq\varnothing$; moreover, $x$ is not universal, so $N_x\cap I\neq I$. Hence the six sets counted by the $\pi_u$ and the $\kappa_u$ partition $K$, which gives (1), and $S$ is recovered, up to isomorphism, from their sizes, which gives (3). Since $|N_u-N_v|=\pi_u+\kappa_v$, (2) follows from \Cref{prop.basicas.sigma_uv}. Finally, if $\sigma_{ac}=0$, then by (2) we may assume, exchanging $a$ and $c$ if necessary, that $\pi_a+\kappa_c=0$, and (4) follows from (1) and (2).
\end{proof}

Following \cite{pastine2025simple}, we say that a split graph $(S,K,I)$ is \emph{$\delta$-homogeneous} if $\deg_S(v)=\delta$ for all $v\in I$. Recall that if $S$ is active, then $|K|=\omega(S)$, $|I|=\alpha(S)$, and $(K,I)$ is the unique bipartition for $S$. Throughout the next sections we use the following facts.

\begin{proposition}[\cite{pastine20252,pastine2025simple}]
	\label{hechos.previos}
	Let $(S,K,I)$ be a split graph and let $\Phi=\Phi(S,K,I)$.
	\begin{enumerate}
		\item If $G$ is a graph, then
		\[ \deg((S,K,I)\circ G)=\deg(S)+\deg(G), \]
		and $S\circ G$ is active if and only if $S$ and $G$ are active.
		
		\item If $K=\bigcup_{v\in I}N_v$, $|I|\geq 2$ and $\Phi$ is simple and connected, then $(S,K,I)$ is $\delta$-homogeneous.
		
		\item If $(S,K,I)$ is $\delta$-homogeneous and $\deg(S)\geq 1$, then:
		\begin{enumerate}
			\item $\sigma_{uv}=0$ if and only if $N_u=N_v$;
			\item $\Phi$ contains no three vertices $u,v,w$ such that $\sigma_{uv}\neq 0$ and $\sigma_{uw}=0=\sigma_{vw}$.
		\end{enumerate}
		%\item If $K=\bigcup_{v\in I}N_v$, $|I|\geq 2$, $\Phi$ is simple and complete, and $S$ is active, then $\omega(S)=\alpha(S)$ and $\delta\in\{1,\omega(S)-1\}$.
	\end{enumerate}
\end{proposition}

%%%%%%%%%%%%%%%%%%%%%%%%%%%%%%%%%%%%%%%%%%%%%%%%%%%%%%%%%%%%%%%%%%%%%%

\subsection{Modules and induced 4-cycles}\label{sec_tools_C4}

A set $M\subseteq V(G)$ is a \emph{module} of $G$ if every vertex of
$V(G)- M$ is adjacent either to all vertices of $M$ or to none of
them (see \cite{gallai1967}). If $M$ is a module of $G$ and $W\subseteq V(G)$, then $M\cap W$ is a
module of $G[W]$.

\begin{lemma}
	\label{modulo.reemplazo}
	Let $M$ be a module of $G$, let $H\prec G$ and put $W=V(H)\cap M$. If
	$A\subseteq M$ and $\varphi\colon W\to A$ is an isomorphism from
	$G[W]$ onto $G[A]$, then $G\bigl[(V(H)-W)\cup A\bigr]\approx H$.
\end{lemma}

\begin{proof}
	Extend $\varphi$ by the identity on $V(H)-W$. Adjacencies inside
	$V(H)-W$ are unchanged and those inside $W$ are preserved by
	$\varphi$. Since $V(H)-W\subseteq V(G)-M$ and $M$ is a module, every
	$v\in V(H)-W$ is adjacent to all or to none of the vertices of $M$;
	hence $uv\in E(G)$ if and only if $\varphi(u)v\in E(G)$, for all
	$u\in W$.
\end{proof}

Let $A_4(G)$ be the graph on $V(G)$ in which two vertices are adjacent if
and only if they lie in a common member of $\mathcal{Q}_G$, that is, if some
2-switch on $G$ involves both of them; $A_4(G)$ is the $2$-section of the
$A_4$-structure of $G$ introduced in \cite{barrus.west.A4}.

\begin{theorem}[\cite{barrus.west.A4}; see also \cite{pastine2025simple}]
	\label{indecomp.characterization}
	If $G_n\circ\cdots\circ G_1$ is the Tyshkevich decomposition of $G$,
	then $A_4(G)=\dot{\bigcup}_{i=1}^{n}A_4(G_i)$. In particular, $G$ is
	indecomposable if and only if $A_4(G)$ is connected.
\end{theorem}

Note that \eqref{deg.como.suma} gives
$\sum_{F\in\mathcal{F}}\deg(F)\leq\deg(G)$ for every
$\mathcal{F}\subseteq\mathcal{Q}_G$, with equality if and only if
$\mathcal{F}=\mathcal{Q}_G$. If $C\in\mathcal{Q}_G(C_4)$ and $v\in V(G)-V(C)$, define
\[
\mathcal{H}_v=\bigl\{H\in \mathcal{Q}_G:\ V(H)\subseteq\{v\}\cup V(C)
\ \text{ and }\ V(H)\neq V(C)\bigr\}.
\]

Let $M\subseteq V(G)$. A vertex $v\in V(G)-M$ \emph{distinguishes} $M$ if $\varnothing\neq N_G(v)\cap M\neq M$; thus $M$ is a module of $G$ if and only if no vertex distinguishes it.

\begin{lemma}
	\label{conteo.local.C4}
	Let $C=c_0c_1c_2c_3c_0$ be an induced $C_4$ in $G$, with subscripts
	in $\mathbb{Z}_4$, and let $v\in V(G)-V(C)$. Then
	\[
	\sum_{H\in\mathcal{H}_v}\deg(H)=
	\begin{cases}
		0, & \text{if } v \text{ does not distinguish } V(C);\\[2pt]
		4, & \text{if } N_G(v)\cap V(C)=\{c_i,c_{i+2}\}
		\text{ for some } i\in\mathbb{Z}_4;\\[2pt]
		2, & \text{otherwise.}
	\end{cases}
	\]
\end{lemma}

\begin{proof}
	The $4$-subsets of $\{v\}\cup V(C)$ other than $V(C)$ are exactly the
	sets $W_i=\{v,c_{i-1},c_i,c_{i+1}\}$ with $i\in\mathbb{Z}_4$, so
	$\mathcal{H}_v=\{G[W_i]:i\in\mathbb{Z}_4\}\cap \mathcal{Q}_G$. Put
		$N=N_G(v)\cap V(C)$. Since $c_{i-1}c_i,c_ic_{i+1}\in E(G)$ and
	$c_{i-1}c_{i+1}\notin E(G)$, checking the eight possibilities for
	$N\cap\{c_{i-1},c_i,c_{i+1}\}$ gives
	\[
	G[W_i]\in \mathcal{Q}_G\iff c_i\notin N\ \text{ and }\
	N\cap\{c_{i-1},c_{i+1}\}\neq\varnothing,
	\]
	and in that case $\deg(G[W_i])=|N\cap\{c_{i-1},c_{i+1}\}|$, since
	$G[W_i]\approx P_4$ when this number is $1$ and $G[W_i]\approx C_4$
	when it is $2$.
	
	If $v$ does not distinguish $V(C)$, then $N\in\{\varnothing,V(C)\}$,
	no $W_i$ satisfies the condition, and the sum is $0$. Otherwise,
	rotating $C$ we may assume $N\in\{\{c_0\},\{c_0,c_1\},\{c_0,c_2\},
	\{c_0,c_1,c_2\}\}$. If $N=\{c_0,c_2\}$, then
	$\mathcal{H}_v=\{G[W_1],G[W_3]\}$ with both graphs isomorphic to
	$C_4$, so the sum is $4$. In each of the other three cases the sum is
	$2$: for $N=\{c_0\}$ we get $\mathcal{H}_v=\{G[W_1],G[W_3]\}$ with both
	graphs isomorphic to $P_4$; for $N=\{c_0,c_1\}$ we get
	$\mathcal{H}_v=\{G[W_2],G[W_3]\}$, again two copies of $P_4$; and for
	$N=\{c_0,c_1,c_2\}$ we get $\mathcal{H}_v=\{G[W_3]\}$ with
	$G[W_3]\approx C_4$.
\end{proof}

\begin{theorem}
	\label{C4.distinguidores}
	Let $C\prec G$ with $C\approx C_4$. Then
	\begin{equation}
		\label{conteo.global.C4}
		\deg(G)\ \geq\ 2+\sum_{v\in V(G)-V(C)}\ \sum_{H\in\mathcal{H}_v}\deg(H),
	\end{equation}
	with equality if and only if
	$\mathcal{Q}_G=\{C\}\cup\bigcup_{v\in V(G)-V(C)}\mathcal{H}_v$. Consequently,
	at most $\lfloor(\deg(G)-2)/2\rfloor$ vertices distinguish $V(C)$. In
	particular, $V(C)$ is a module of $G$ if $\deg(G)\leq 3$; and if
	$\deg(G)\leq 5$, then at most one vertex $v$ distinguishes $V(C)$,
	and $N_G(v)\cap V(C)$ is not a pair of opposite vertices of $C$.
\end{theorem}

\begin{proof}
	Every member of $\mathcal{H}_v$ contains $v$ and exactly three
	vertices of $C$; hence the families $\mathcal{H}_v$ are pairwise
	disjoint and none of them contains $C$. Since $\deg(C)=2$,
	\eqref{conteo.global.C4} and its equality case follow from
	\eqref{deg.como.suma}.
	
	By \Cref{conteo.local.C4}, a vertex contributes $0$ to the sum in
	\eqref{conteo.global.C4} if it does not distinguish $V(C)$, and at
	least $2$ if it does. If $t$ vertices distinguish $V(C)$, then
	\eqref{conteo.global.C4} gives $\deg(G)\geq 2+2t$, so
	$t\leq\lfloor(\deg(G)-2)/2\rfloor$. For $\deg(G)\leq 3$ this yields
	$t=0$, that is, $V(C)$ is a module; for $\deg(G)\leq 5$ it yields
	$t\leq 1$. Finally, if $\deg(G)\leq 5$ and a vertex $v$ had
	$N_G(v)\cap V(C)$ equal to a pair of opposite vertices of $C$, then
	its contribution would be $4$ by \Cref{conteo.local.C4}, whence
	$\deg(G)\geq 2+4=6$, a contradiction.
\end{proof}

\begin{lemma}
	\label{C4.modulo.deg>=6}
	Let $G$ be a prime graph and let $C\prec G$ with $C\approx C_4$. If
	$V(C)$ is a module of $G$ and $G\neq C$, then $\deg(G)\geq 6$.
\end{lemma}

\begin{proof}
	Put $M=V(C)$. Since $G\neq C$, we have $V(G)-M\neq\varnothing$, and
	since $A_4(G)$ is connected by \Cref{indecomp.characterization}, some
	$H\in \mathcal{Q}_G$ meets both $M$ and $V(G)-M$. Put $W=V(H)\cap M$, so
	that $1\leq|W|\leq 3$. If $|W|=3$, the vertex of $V(H)-W$ is adjacent
	to all or to none of the vertices of $W$, which is impossible because the vertices of
	$P_4$, $C_4$ and $2K_2$ have degree 1 or 2. If $|W|=2$, say $W=\{w_1,w_2\}$, then since $M$ is a module
	every vertex of $V(H)-W$ is adjacent to both or to neither of $w_1,w_2$;
	as no two vertices of $P_4$ have this property, $\deg(H)=2$. By \Cref{modulo.reemplazo},
	each $A\subseteq M$ with $G[A]\approx G[W]$ yields a member
	$G[(V(H)-W)\cup A]$ of $\mathcal{Q}_G$ isomorphic to $H$ and distinct from
	$C$; since $C_4$ has four pairs of vertices inducing an edge and two
	inducing a non-edge, there are at least two such $A$, whence
	$\deg(G)\geq 2+2\cdot 2=6$ by \eqref{deg.como.suma}. If $|W|=1$,
	\Cref{modulo.reemplazo} applied to the four singletons
	$A\subseteq M$ yields four members of $\mathcal{Q}_G$ distinct from $C$, so
	$\deg(G)\geq 2+4=6$.
\end{proof}


%%%%%%%%%%%%%%%%%%%%%%%%%%%%%%%%%%%%%%%%%%%%%%%%%%%%%%%%%%%%%%

\section{Active graphs of degree 1 and 2}\label{sec_deg<3}

An immediate consequence of  \Cref{|G^*|} is that $P_4$ is the only active graph of degree 1. Then, by \Cref{hechos.previos}, $P_4^2$ is the only decomposable active graph of degree 2.
Now, consider a prime graph $G$ such that $\deg(G)=|\mathcal{Q}_G(P_4)|=2$. By  \Cref{5>degG=|{P_4}|.implica.G.split}, $G$ is split. 
%Let $n=|G|$. Applying  \Cref{|G^*|}, we obtain $4 \leq n \leq 8$. However, it is clear that $n \neq 4$, since otherwise we would have $|\mathcal{Q}_G(P_4)|=1$. 
%Assume $G=S \circ H$ for some graphs $S$ and $H$ with $|S|,|H|\geq 1$. Since $G$ is active, it follows from  \Cref{SoG.act.iff.S_G.act} that $S$ and $H$ are active. So, $\deg(S)=1=\deg(H)$ by \Cref{hechos.previos}. Consequently, $S \approx P_4\approx H$, and hence $G \approx P_4^2$ when $G$ is decomposable. 

Then $\Phi(G)$ is connected, by  \Cref{S.primo.iff.Phi(S).conexo}. This means that $\Phi(G)$ is isomorphic to $K_2$ or $P_3$, ignoring multiplicities. If $(K,I)$ is the bipartition for $G$, we have $(|K|,|I|)=(\omega,\alpha)$, where $\omega=\omega(G)$ and $\alpha=\alpha(G)$. Assume $\alpha\leq\omega$. As $5 \leq \alpha+\omega=|G|$, it follows that $\omega \neq 2$. Thus, $(\omega,\alpha)\in \{(3,2), (3,3)\}$ by \Cref{|I|_|K|<=2deg(S)}.
If $\alpha=2$, then $G\approx D_{2,1}$ by \Cref{S_activo_|I|=2} (see \Cref{grafos.activos.deg<3}).
If $\alpha=3$, then $\Phi=abc$, and \Cref{S_activo_|I|=3} gives $\kappa_b(\pi_b+\kappa_a)=1=\pi_b(\pi_c+\kappa_b)$. Hence $\pi_b=\kappa_b=1$ and $\kappa_a=\pi_c=0$, so $\omega=2$, which is absurd.
%If $\alpha=3$, then $\Phi=abc$, and so $d_a=d_b=d_c$ and $N_a=N_c$, by \Cref{prop.basicas.sigma_uv}. Since $G$ does not have universal vertices, $\varnothing=(N_a\cap N_c)\cap N_b=N_a\cap N_b$. Then, $1=\sigma_{ab}=d_a d_b$, by \Cref{prop.basicas.sigma_uv}. Since $K=N_a\cup N_b$, we obtain $3=\omega=d_a+d_b-\eta_{ab}=2$, which is absurd.
%If $\alpha=2$, let $I=\{a,b\}$, so that $\sigma_{ab}=2$. Clearly, $\eta_{ab}=0$, since $G$ has no universal vertices. Thus, $3=\omega=d_a+d_b$, and hence $(d_a,d_b)=(1,2)$. Up to isomorphism, there is a unique split graph with these characteristics; we call it $D_{2,1}$ and it can be seen in  \Cref{grafos.activos.deg<3}.

%\begin{theorem}[\cite{pastine2025simple}]
%	\label{caract.Phi.simple.conexo}
%	Given a split graph $(S,K,I)$ with $2\leq|K|\leq |I|$ and $K=\bigcup_{v \in I}N_S(v)$, consider the associated split graph $(R,K,I)$ such that $\deg_R(v)=1$ for all $v\in I$, where $u\cong_R v$ if and only if $\sigma_{uv}(S)=0$, for all $\{u,v\}\subseteq I$. Then $S$ is active and $\Phi(S)$ is simple and connected, if and only if $S$ is isomorphic to $R$ or $\overline{R^{\iota}}$.
%\end{theorem}


%we deduce from  \Cref{caract.Phi.simple.conexo} that $G$ is isomorphic to $(R,[3],I)$ or $\overline{R^{\iota}}$, where $N_R(a)=N_R(c)=\{1\}$ and $N_R(b)=\{2\}$. This is impossible because $G$ is active, but neither $R$ nor $\overline{R^{\iota}}$ are (vertex 3 is swing in $R$ and universal in $\overline{R^{\iota}}$). \ap{revisar!}

So, if $G$ is a prime split graph of degree 2, then $G\approx D_{2,1}$ or $\overline{D_{2,1}}$. If $G$ is a prime non-split graph of degree 2, it is clear that $|\mathcal{Q}_G|=1$. Moreover, $|G|=4$, since otherwise $|\mathcal{Q}_G|\geq 2$. Then, $G\approx C_4$ or $2K_2$.
%\begin{equation}
%	\label{eq35}
%	\mathcal{Q}_G=\mathcal{Q}_G(C_4) \cup \mathcal{Q}_G(2K_2).
%\end{equation}
%Recall that $4 \leq n \leq 8$, by  \Cref{|G^*|}. Obviously, $G \approx C_4$ or $2K_2$ if $n=4$. On the other hand, if $n \geq 5$, then $|\mathcal{Q}_G| \geq 2$. This would imply $\deg(G) \geq 4$, by \eqref{eq35}, contradicting that $\deg(G)=2$. Clearly, $C_4$ and $2K_2$ are indecomposable. Indeed, if there were graphs $S$ and $H$ with $|S|,|H|\geq 1$ such that $S \circ H \approx C_4$, both factors would have to be active, by  \Cref{SoG.act.iff.S_G.act}. But then applying  \Cref{|G^*|} leads to the following absurdity: $4=|C_4|=|S|+|H| \geq 8$. The indecomposability of $2K_2$ now follows from $2K_2=\overline{C_4}$.

\begin{theorem}
	\label{clasificacion.activo_deg=2}
	Let $G$ be an active graph.
	\begin{enumerate}
		\item If $\deg(G)=1$, then $G \approx P_4$.
		\item If $\deg(G)=2$ and $G$ is decomposable, then $G \approx P_4^2$.
		\item If $\deg(G)=2$ and $G$ is indecomposable, then $G$ is isomorphic to one of these graphs: $D_{2,1}$, $\overline{D_{2,1}}$, $C_4$, $2K_2$.
	\end{enumerate}
	All graphs mentioned here are represented in  \Cref{grafos.activos.deg<3}.
\end{theorem}

%\begin{figure}[h]
%	\centering
%	\includegraphics[scale=0.8]{grafos_activos_de_grado_1_o_2.pdf}
%	\caption{The 6 active graphs of degree 1 or 2.}
%	\label{grafos.activos.deg<3}
%\end{figure}
\begin{proof}
	It follows from the previous discussion.
\end{proof}

\begin{figure}[h!]
	\centering
	\begin{tikzpicture}[
		scale=1,
		every node/.style={circle, draw, inner sep=2pt},
		lbl/.style={draw=none}
		]
		
		% ===================== D5 =====================
		\begin{scope}[shift={(0,2)}]
			\node (t) at (0,1.2) {};
			\node (l) at (-0.6,0.4) {};
			\node (r) at (0.6,0.4) {};
			\node (b) at (0,-0.4) {};
			\node (e) at (0.6,-0.4) {};
			
			\draw (t)--(l)--(b)--(r)--(t);
			\draw (l)--(r);
			\draw (b)--(e);
			
			\node[lbl] at (1.2,-0.4) {$D_{2,1}$};
		\end{scope}
		
		% ===================== complement D5 =====================
		\begin{scope}[shift={(2.4,2)}]
			\node (t) at (0,1.2) {};
			\node (l) at (-0.6,0.4) {};
			\node (r) at (0.6,0.4) {};
			\node (b) at (0,-0.4) {};
			\node (e) at (0.6,-0.4) {};
			
			\draw (t)--(r);
			\draw (l)--(r);
			\draw (r)--(b);
			\draw (b)--(e);
			
			\node[lbl] at (1.2,0.4) {$\overline{D_{2,1}}$};
		\end{scope}
		
		% ===================== P4^2 =====================
		\begin{scope}[shift={(4.8,2)}]
			
			% fila superior
			\node (t1) at (0,1) {};
			\node (t2) at (1,1) {};
			\node (t3) at (2,1) {};
			\node (t4) at (3,1) {};
			
			% fila inferior
			\node (b1) at (0,0) {};
			\node (b2) at (1,0) {};
			\node (b3) at (2,0) {};
			\node (b4) at (3,0) {};
			
			% caminos
			\draw (t1)--(t2)--(t3)--(t4);
			\draw (b1)--(b2)--(b3)--(b4);
			
			% verticales (solo las dos centrales)
			\draw (t2)--(b2);
			\draw (t3)--(b3);
			
			% diagonales EXACTAS del dibujo
			\draw (b1)--(t2);
			\draw (b1)--(t3);
			\draw (b2)--(t3);
			\draw (b4)--(t2);
			\draw (b3)--(t2);
			\draw (b4)--(t3);
			
			\node[lbl] at (1.5,-0.7) {$P_4^2$};
		\end{scope}
		
		% ===================== C4 =====================
		\begin{scope}[shift={(0,-1)}]
			\node (a) at (0,1) {};
			\node (b) at (1,1) {};
			\node (c) at (1,0) {};
			\node (d) at (0,0) {};
			
			\draw (a)--(b)--(c)--(d)--(a);
			
			\node[lbl] at (0.5,1.4) {$C_4$};
		\end{scope}
		
		% ===================== 2K2 =====================
		\begin{scope}[shift={(2.4,-1)}]
			\node (a) at (0,1) {};
			\node (b) at (1,1) {};
			\node (c) at (0,0) {};
			\node (d) at (1,0) {};
			
			\draw (a)--(b);
			\draw (c)--(d);
			
			\node[lbl] at (0.5,1.4) {$2K_2$};
		\end{scope}
		
		% ===================== P4 =====================
		\begin{scope}[shift={(4.8,-1)}]
			\node (a) at (0,0) {};
			\node (b) at (1,0) {};
			\node (c) at (2,0) {};
			\node (d) at (3,0) {};
			
			\draw (a)--(b)--(c)--(d);
			
			\node[lbl] at (1.5,0.6) {$P_4$};
		\end{scope}
		
	\end{tikzpicture}
	\caption{The 6 active graphs of degree 1 or 2.}
	\label{grafos.activos.deg<3}
\end{figure}



%%%%%%%%%%%%%%%%%%%%%%%%%%%%%%%%%%%%%%%%%%%%%%%%%%%%%%%%%%%

\section{Active graphs of degree 3}\label{sec_deg=3}

Consider an active graph $G$ such that $\deg(G)=|\mathcal{Q}_G(P_4)|=3$. Thanks to  \Cref{5>degG=|{P_4}|.implica.G.split}, we know that $G$ is split. Suppose $G$ is decomposable. Since $G$ is active, it follows from \Cref{hechos.previos} that each factor of $G$ is an active split graph of degree 1 or 2. Consequently,
\[ G \in \{P_4 \circ D_{2,1}, D_{2,1} \circ P_4, P_4 \circ \overline{D_{2,1}}, \overline{D_{2,1}} \circ P_4, P_4^3\}, \]
by  \Cref{clasificacion.activo_deg=2}.

If instead $(G,K,I)$ is indecomposable, assume $\alpha=|I|\leq|K|=\omega$. Clearly, $\omega \geq 3$, since otherwise we would have $G\approx P_4$. Thus, $2 \leq \alpha \leq 4$ and $\omega \in \{3,4\}$, by \Cref{|I|_|K|<=2deg(S)}. Applying \Cref{S.primo.iff.Phi(S).conexo}, we have
\[ \Phi(G) \in \{K_2, K_3, P_3, K_{3,1}, P_4\}, \]
ignoring multiplicities if any (see \Cref{Phi.posibles.deg=3}).
\begin{figure}[h]
	\centering
	\begin{tikzpicture}[scale=1, every node/.style={draw, circle, fill=white, minimum size=0.25cm, inner sep=1pt}]
		
		% triple edge
		\begin{scope}[shift={(0,0)}]
			\node (a1) at (0,0) {};
			\node (a2) at (0,1) {};
			\draw (a1) -- (a2);
			\draw[bend left=25] (a1) to (a2);
			\draw[bend right=25] (a1) to (a2);
		\end{scope}
		
		% double edge with a pendant vertex
		\begin{scope}[shift={(1.7,0)}]
			\node (b1) at (0,0) {};
			\node (b2) at (1,0) {};
			\node (b3) at (0,1) {};
			\draw (b1) -- (b2);
			\draw (b1) -- (b3);
			\draw[bend right=25] (b1) to (b3);
		\end{scope}
		
		% K_3
		\begin{scope}[shift={(4.3,0)}]
			\node (c1) at (0,0) {};
			\node (c2) at (1,0) {};
			\node (c3) at (0,1) {};
			\draw (c1) -- (c2);
			\draw (c1) -- (c3);
			\draw (c2) -- (c3);
		\end{scope}
		
		% P_4
		\begin{scope}[shift={(6.9,0)}]
			\node (d1) at (0,0) {};
			\node (d2) at (1,0) {};
			\node (d3) at (0,1) {};
			\node (d4) at (1,1) {};
			\draw (d1) -- (d2);
			\draw (d1) -- (d3);
			\draw (d2) -- (d4);
		\end{scope}
		
		% S_4
		\begin{scope}[shift={(9.5,0)}]
			\node (e0) at (0.5,0.5) {};
			\node (e1) at (0,1) {};
			\node (e2) at (1,1) {};
			\node (e3) at (0.5,0) {};
			\draw (e0) -- (e1);
			\draw (e0) -- (e2);
			\draw (e0) -- (e3);
		\end{scope}
		
	\end{tikzpicture}
	\caption{The 5 connected multigraphs of size 3.}
	\label{Phi.posibles.deg=3}
\end{figure}
%We now apply  \Cref{prop.basicas.sigma_uv} to this class of split graphs, obtaining substantial understanding of the structure of both $S$ and $\Phi(S)$.
If $\alpha=2$, then $G\approx D_{3,1}$ by \Cref{S_activo_|I|=2} (see \Cref{grafos.primos.deg=3}).
If $\Phi=\Phi(G) \in \{K_3, K_{3,1}, P_4\}$, then $G$ is $\delta$-homogeneous, by  \Cref{hechos.previos}.  The same proposition therefore prohibits $\Phi=P_4$. Suppose $\Phi=K_{3,1}$. Thanks to  \Cref{hechos.previos}, we know that $N_1=N_2=N_3$. On the other hand, $0=|\bigcap_{i=1}^4 N_i|=\eta_{14}$, because $G$ has no universal vertices. This means that $\delta=1$, by  \Cref{prop.basicas.sigma_uv}. But then $\omega=2$, contradicting $\alpha \leq \omega$. Hence, $\alpha\neq 4$.

If $\alpha=3$, then $\Phi=abca$, or $\Phi=abc$ with $\{\sigma_{ab},\sigma_{bc}\}=\{1,2\}$. If $\Phi=abca$ then the six factors in \Cref{S_activo_|I|=3}(2) are equal to $1$, so $\pi_a=\pi_b=\pi_c$, $\kappa_a=\kappa_b=\kappa_c$ and $\pi_a+\kappa_a=1$. This yields two graphs, complementary to each other; we call $T_{111}$ the one with $\pi_a=1$ (the subscript records the degrees of the $I$-vertices). If $\Phi=abc$, then, exchanging $a$ and $c$ if necessary, \Cref{S_activo_|I|=3}(4) gives $\sigma_{ab}=\kappa_b(\pi_b+\kappa_a)$ and $\sigma_{bc}=\pi_b(\pi_c+\kappa_b)$, so $\pi_b=\kappa_b=1$. If $\sigma_{ab}=2$, then $\kappa_a=1$ and $\pi_c=0$, and we call the resulting graph $T_{221}$; if $\sigma_{ab}=1$, then $\kappa_a=0$ and $\pi_c=1$, which yields $\overline{T_{221}}$. In all cases $\omega=3$. These graphs are shown in \Cref{grafos.primos.deg=3}.
%If $\alpha=3$, then $\Phi=K_3$ or $\Phi=P_3$ with a double-edge. If $\Phi=K_3$, \Cref{hechos.previos} tells us that $\omega=3$ and $\delta \in \{1,2\}$. Both possibilities for $\delta$ are realized, forming a pair of graphs that are complements of each other. We name $T_{111}$ the graph resulting from taking $\delta=1$ (the subscript records the degrees of the $I$-vertices). These are shown in  \Cref{grafos.primos.deg=3}. If $\Phi=abc$, with $\sigma_{ab}=2$ and $\sigma_{bc}=1$, then two cases arise: 1) $d_a < d_b$; 2) $d_a > d_b$. Suppose $d_a < d_b$. By \Cref{prop.basicas.sigma_uv}, we have $d_c = d_b = d_a + 1$ and $N_a \subseteq N_c$. Thus, $\eta_{ac}=d_a$. Since $G$ contains no universal vertices, it follows that $\varnothing = (N_a \cap N_c) \cap N_b = N_a \cap N_b$, and so $d_a=1$, by \Cref{prop.basicas.sigma_uv}. Since $1=\sigma_{bc}=((d_a+1)-\eta_{bc})^2$, we have $\eta_{bc}=d_a=1$. Then, applying inclusion-exclusion principle to $K=N_a\cup N_b\cup N_c$, we find that $\omega=3$. All these characteristics determine a graph we name $T_{221}$, which can be seen in  \Cref{grafos.primos.deg=3}. If $d_a > d_b$, an almost identical reasoning leads us to $\overline{T_{221}}$.


%If $\alpha=2$, let $I=\{a,b\}$, so that $\sigma_{ab}=3$. Without loss of generality, we can assume $d_a \leq d_b$. Since $G$ has no universal vertices, we have $\eta_{ab}=0$. Thus, $d_a=1$ and $d_b=3$ by \Cref{prop.basicas.sigma_uv}, which implies $\omega=4$. Up to isomorphism, there is a unique split graph with these characteristics; we name it $D_{3,1}$ and it can be seen in \Cref{grafos.primos.deg=3}.\\
\begin{figure}[ht]
	\centering
	\begin{tikzpicture}[scale=0.9, 
		every node/.style={draw, circle, fill=white, inner sep=0pt, minimum size=6.5pt},
		]
		
		%%%%%% T_{111} %%%%%%
		\begin{scope}[shift={(0,3.5)}]
			\node (t1) at (0, 1) {};
			\node (t2) at (-0.55, 0) {};
			\node (t3) at (0.55, 0) {};
			\node (p1) at (0, 2) {};
			\node (p2) at (-1.25, -0.7) {};
			\node (p3) at (1.25, -0.7) {};
			\draw (t1) -- (t2) -- (t3) -- (t1);
			\draw (t1) -- (p1);
			\draw (t2) -- (p2);
			\draw (t3) -- (p3);
			\node[draw=none] at (-1, 1) {$T_{111}$};
		\end{scope}
		
		%%%%%% D_{3,1} %%%%%%
		\begin{scope}[shift={(2.2,3)}]
			\node (a)  at (0, 2.2) {};
			\node (k1) at (1.0, 1.3) {};
			\node (k2) at (2.2, 1.3) {};
			\node (k3) at (1.0, 0) {};
			\node (k4) at (2.2, 0) {};
			\node (pp) at (3.2, 0) {};
			\draw (k1) -- (k2) -- (k4) -- (k3) -- (k1);
			\draw (k1) -- (k4);
			\draw (k2) -- (k3);
			\draw (a) -- (k1);
			\draw (a) -- (k2);
			\draw (a) -- (k3);
			\draw (k4) -- (pp);
			\node[draw=none] at (2, 2) {$D_{3,1}$};
		\end{scope}
		
		%%%%%% T_{221} %%%%%%
		\begin{scope}[shift={(6.2,3.5)}]
			\node (u1) at (0, 1.3) {};
			\node (u2) at (1, 1.3) {};
			\node (u3) at (2, 1.3) {};
			\node (u4) at (0.5, 0.3) {};
			\node (u5) at (1.5, 0.3) {};
			\node (u6) at (1, -0.7) {};
			\draw (u1) -- (u2) -- (u3);
			\draw (u1) -- (u4);
			\draw (u3) -- (u5);
			\draw (u2) -- (u4);
			\draw (u2) -- (u5);
			\draw (u4) -- (u5);
			\draw (u4) -- (u6);
			\node[draw=none] at (2, -0.7) {$T_{221}$};
		\end{scope}
		
		%%%%%% \overline{T_{111}} %%%%%%
		\begin{scope}[shift={(0,0.5)}]
			\node (b1) at (-1.2, 1) {};
			\node (b2) at (0, 1) {};
			\node (b3) at (1.2, 1) {};
			\node (b4) at (-0.55, 0) {};
			\node (b5) at (0.55, 0) {};
			\node (b6) at (0, -1.1) {};
			\draw (b1) -- (b2) -- (b3);
			\draw (b1) -- (b4);
			\draw (b3) -- (b5);
			\draw (b2) -- (b4);
			\draw (b2) -- (b5);
			\draw (b4) -- (b5);
			\draw (b4) -- (b6);
			\draw (b5) -- (b6);
			\node[draw=none] at (-1, -1) {$\overline{T_{111}}$};
		\end{scope}
		
		%%%%%% \overline{D_{3,1}} %%%%%%
		\begin{scope}[shift={(2.6,0)}]
			\node (c1) at (0, 1.5) {};
			\node (c2) at (0.9, 0.6) {};
			\node (c3) at (2.0, 0.6) {};
			\node (c4) at (0.9, -0.6) {};
			\node (c5) at (2.0, -0.6) {};
			\node (c6) at (3.1, -0.6) {};
			\draw (c1) -- (c2);
			\draw (c2) -- (c3);
			\draw (c2) -- (c4);
			\draw (c2) -- (c5);
			\draw (c5) -- (c6);
			\node[draw=none] at (-0.1, 0) {$\overline{D_{3,1}}$};
		\end{scope}
		
		%%%%%% \overline{T_{221}} %%%%%%
		\begin{scope}[shift={(6.2,0)}]
			\node (d1) at (0, 1.3) {};
			\node (d2) at (1, 1.3) {};
			\node (d3) at (2, 1.3) {};
			\node (d4) at (0.5, 0.3) {};
			\node (d5) at (1.5, 0.3) {};
			\node (d6) at (1, -0.7) {};
			\draw (d2) -- (d3);
			\draw (d1) -- (d4);
			\draw (d2) -- (d4);
			\draw (d2) -- (d5);
			\draw (d4) -- (d5);
			\draw (d4) -- (d6);
			\draw (d5) -- (d6);
			\node[draw=none] at (2, -0.7) {$\overline{T_{221}}$};
		\end{scope}
		
	\end{tikzpicture}
	\caption{The 6 prime graphs of degree $3$.}
	\label{grafos.primos.deg=3}
\end{figure}

So far, we have classified active graphs $G$ of degree 3 such that $\mathcal{Q}_G=\mathcal{Q}_G(P_4)$. To complete the classification of all active graphs of degree 3, it remains to analyze the following case:
\begin{equation}
	\label{eq38}
	|\mathcal{Q}_G(C_4)| + |\mathcal{Q}_G(2K_2)| = 1 = |\mathcal{Q}_G(P_4)|.
\end{equation}
Recall that, by definition of $\circ$, if $G=X \circ Y$, then $X$ is split. So, if $G$ is decomposable and satisfies \eqref{eq38}, then $G \in \{P_4 \circ C_4, P_4 \circ 2K_2\}$, by \Cref{clasificacion.activo_deg=2} and \Cref{hechos.previos}. In the following lemma, we prove that there are no prime graphs satisfying \eqref{eq38}.





%\begin{lemma}
%	\label{C4.es.modulo}
%	Let $G$ be a graph with $|\mathcal{Q}_G|\in\{1,2\}$ and  $\mathcal{Q}_G(C_4)=\{C\}$. Then $V(C)$ is a module of $G$.
%\end{lemma}
%
%\begin{proof}
%	Write $C=c_0c_1c_2c_3c_0$, let $v\in V(G)- V(C)$, put
%	$N=N_v\cap V(C)$ and, for $i\in\mathbb{Z}_4$, put
%	$W_i=\{v,c_{i-1},c_i,c_{i+1}\}$. Since $c_{i-1}c_i,c_ic_{i+1}\in E(G)$ and
%	$c_{i-1}c_{i+1}\notin E(G)$, checking the eight possibilities for
%	$N_v\cap\{c_{i-1},c_i,c_{i+1}\}$ gives
%	\[
%	G[W_i]\in \mathcal{Q}_G\iff c_i\notin N\ \text{ and }\ N\cap\{c_{i-1},c_{i+1}\}
%	\neq\varnothing,
%	\]
%	and in that case $G[W_i]\approx P_4$ if $|N\cap\{c_{i-1},c_{i+1}\}|=1$,
%	while $G[W_i]\approx C_4$ if $|N\cap\{c_{i-1},c_{i+1}\}|=2$.
%	
%	Suppose $\varnothing\neq N\neq V(C)$. Rotating $C$ we may assume
%	$N\in\{\{c_0\},\{c_0,c_1\},$ $\{c_0,c_2\},\{c_0,c_1,c_2\}\}$. In the first
%	three cases exactly two of the graphs $G[W_i]$ lie in $\mathcal{Q}_G$, and all of them
%	contain $v$, so $|\mathcal{Q}_G|\geq 3$. In the last case $G[W_3]\approx C_4$, so $|\mathcal{Q}_G(C_4)|\geq 2$. Both conclusions are absurd.
%	Hence $N\in\{\varnothing,V(C)\}$.
%\end{proof}

\begin{lemma}
	\label{no.existe.deg3.primo}
	There exists no prime graph satisfying \eqref{eq38}.
\end{lemma}

\begin{proof}
	Assume by contradiction that $G$ is prime and satisfies \eqref{eq38}. A
	$4$-subset of $V(G)$ induces $P_4$, $C_4$ or $2K_2$ in $G$ if and only if
	it induces $P_4$, $2K_2$ or $C_4$ in $\overline{G}$, respectively, and $G$
	is prime if and only if $\overline{G}$ is; hence we may assume
	$\mathcal{Q}_G=\{P,C\}$ with $P\approx P_4$ and $C\approx C_4$. By
	\eqref{deg.como.suma}, $\deg(G)=3$, so $V(C)$ is a module
	of $G$ by \Cref{C4.distinguidores}. Moreover $G\neq C$, because
	$\deg(C)=2\neq\deg(G)$. Therefore $\deg(G)\geq 6$ by
	\Cref{C4.modulo.deg>=6}, a contradiction.
\end{proof}

%\begin{lemma}
%	\label{no.existe.deg3.primo.5_8}
%	There exists no prime graph $G$ of order 5 or 8 such that
%	\begin{equation}
%		\label{eq39}
%		|\mathcal{Q}_G(C_4)| = 1 = |\mathcal{Q}_G(P_4)|, \ \mathcal{Q}_G(2K_2) = \varnothing.
%	\end{equation}
%\end{lemma}
%
%\begin{proof}
%	Suppose there exists a prime graph $G$ of order 5 or 8 satisfying \eqref{eq39}. Then, $\mathcal{Q}_G^*=\{P,C\}$, where $P \approx P_4$, $C \approx C_4$ and $V(P) \cup V(C) = V(G)$. 
%	%Moreover, $A_4(G)$ is connected by \Cref{indecomp.characterization}.
%	%and that both $V(P)$ and $V(C)$ are 4-cliques in $A_4(G)$.
%	
%	If $|G|=8$, then $V(P)\cap V(C)=\varnothing$, so no edge of $A_4(G)$ joins $V(P)$ to $V(C)$. Hence $A_4(G)$ is disconnected and $G$ is decomposable by \Cref{indecomp.characterization}, a contradiction.
%	
%	If $|G|=5$, then $|V(P)\cap V(C)|=3$; write $C=bx_1x_2x_3b$, so that
%	$V(P)=\{a,x_1,x_2,x_3\}$. Since $x_1x_2x_3\prec G$ and $P$ is an induced $P_4$, the vertex $a$ is adjacent to exactly one end of $x_1x_2x_3$, and interchanging $x_1$ with $x_3$ if necessary we may assume
%	$P=ax_1x_2x_3$. If $ab\notin E(G)$, then $ax_1bx_3$ is an induced $P_4$
%	different from $P$, contradicting \eqref{eq39}; hence $ab\in E(G)$. If
%	moreover $ax_3\notin E(G)$ and $bx_2\notin E(G)$, then $abx_3x_2$ is an
%	induced $P_4$ different from $P$, again a contradiction. Thus
%	$ax_3\in E(G)$ or $bx_2\in E(G)$, contradicting
%	$P,C\prec G$.
%	%If $|G|=5$, then $|V(P) \cap V(C)|=3$. Thus, $P=ax_1x_2x_3$ and $C=bx_1x_2x_3b$. Now observe that necessarily $ab \in E(G)$. Otherwise, $ax_1bx_3 \prec G$, which would contradict the hypothesis about $\mathcal{Q}_G^*$. Consequently, $ax_3 \in E(G)$ or $bx_2 \in E(G)$, since otherwise we would have $abx_3x_2 \prec G$. But this contradicts that $P,C \prec G$.
%\end{proof}
%
%\begin{lemma}
%	\label{no.existe.deg3.primo.7}
%	There exist no prime graphs of order 7 satisfying \eqref{eq39}.
%\end{lemma}
%
%\begin{proof}
%	Suppose there exists a prime graph $G$ of order 7 satisfying \eqref{eq39}. Then, $\mathcal{Q}_G^*=\{P,C\}$, where $P \approx P_4$, $C \approx C_4$ and $V(P) \cup V(C) = V(G)$. Since $|G|=7$, it follows that $P$ and $C$ share exactly one vertex, call it $x$. Therefore, we essentially have two cases to analyze: 1) $P=xa_1a_2a_3$; 2) $P=a_1xa_2a_3$. Let $C=xb_1b_2b_3x$.
%	\begin{enumerate}[(1).]
%		\item Observe that $a_1b_1, a_1b_3 \in E(G)$ (otherwise, $\{b_1xa_1a_2, b_3xa_1a_2\} \subseteq \mathcal{Q}_G^*$). Consequently, also $a_1b_2 \in E(G)$ (otherwise, $a_1b_1b_2b_3a_1 \in \mathcal{Q}_G^*$). This in turn implies the presence of $a_2b_1$ in $G$ (otherwise, $b_1a_1a_2a_3 \in \mathcal{Q}_G^*$). Since the path $a_2b_1xb_3$ cannot be induced in $G$, at least one edge between $b_1b_3$ and $a_2x$ is forced to be in $G$. However, this contradicts that $P,C \prec G$.
%		\item Note that $a_2b_1, a_2b_3 \in E(G)$ (otherwise, $\{b_3xa_2a_3, a_2xb_1b_2\} \subseteq \mathcal{Q}_G^*$). Then, also $a_2b_2 \in E(G)$ (otherwise, $a_2b_1b_2b_3a_2 \in \mathcal{Q}_G^*$). Since the path $a_1xa_2b_2$ cannot be induced in $G$, at least one edge between $a_1a_2$ and $b_2x$ is forced to be in $G$. However, this contradicts that $P,C \prec G$.
%	\end{enumerate}
%\end{proof}
%
%\begin{lemma}
%	\label{no.existe.deg3.primo.6}
%	There exist no prime graphs of order 6 satisfying \eqref{eq39}.
%\end{lemma}
%
%\begin{proof}
%	Suppose there exists a prime graph $G$ of order 6 satisfying \eqref{eq39}. Then, $\mathcal{Q}_G^*=\{P,C\}$, where $P \approx P_4$, $C \approx C_4$ and $V(P) \cup V(C) = V(G)$. Since $|G|=6$, it follows that $V(P) \cap V(C) = \{x_1,x_2\}$. We see that $x_1$ and $x_2$ can be positioned in $P$ essentially in 4 ways: 1) $a_1x_1x_2a_2$; 2) $a_1a_2x_1x_2$; 3) $a_1x_1a_2x_2$; 4) $x_1a_1a_2x_2$.
%	\begin{enumerate}[(1).]
%		\item Let $C=b_1b_2x_2x_1b_1$. Note that necessarily $a_1x_2 \in G$ or $b_2x_1 \in G$ (otherwise, $a_1x_1x_2b_2 \prec G$). This contradicts that $P,C \prec G$.
%		\item Let $C=b_1b_2x_2x_1b_1$. Observe that necessarily $a_2x_2 \in G$ or $b_2x_1 \in G$ (otherwise, $a_2x_1x_2b_2 \prec G$). This contradicts that $P,C \prec G$.
%		\item Let $C=b_1x_1b_2x_2b_1$. Note that necessarily $a_1b_1, a_1b_2 \in G$ (otherwise, $\{a_1x_1b_1x_2, a_1x_1b_2x_2\} \subseteq \mathcal{Q}_G^*$). This forces $b_1b_2$ or $a_1x_2$ to be in $G$ (otherwise, $a_1b_1x_2b_2a_1 \prec G$), which contradicts that $P,C \prec G$.
%		\item Let $C=b_1x_1b_2x_2b_1$. Observe that necessarily $a_2b_1, a_2b_2 \in G$ (otherwise, $\{a_1a_2x_2b_1, a_1a_2x_2b_2\} \subseteq \mathcal{Q}_G^*$). This forces $b_1b_2$ or $a_2x_1$ to be edges of $G$ (otherwise, $a_2b_2x_1b_1a_2 \prec G$), which contradicts that $P,C \prec G$.
%	\end{enumerate}
%\end{proof}
%
%\begin{lemma}
%	\label{no.existe.deg3.primo}
%	There exists no prime graph satisfying \eqref{eq38}.
%\end{lemma}
%
%\begin{proof}
%	Suppose $G$ is a prime graph satisfying \eqref{eq38}. Since $|\mathcal{Q}_G^*|=2$, it is clear that $5 \leq |G| \leq 8$, since $G$ is active and $\deg(G)>1$. We have two cases to analyze: 1) $|\mathcal{Q}_G(C_4)|=1$; 2) $\mathcal{Q}_G(C_4)=\varnothing$.
%	\begin{enumerate}[(1).]
%		\item Since $G$ satisfies \eqref{eq39}, \Cref{no.existe.deg3.primo.5_8,no.existe.deg3.primo.7,no.existe.deg3.primo.6} are contradicted.
%		\item Since $\overline{2K_2}=C_4$ and $\overline{P_4}=P_4$, it follows that $\overline{G}$ satisfies \eqref{eq39}. This contradicts \Cref{no.existe.deg3.primo.5_8,no.existe.deg3.primo.7,no.existe.deg3.primo.6}.
%	\end{enumerate}
%\end{proof}

\begin{theorem}
	\label{clasificacion.activo_deg=3}
	Let $G$ be an active graph of degree 3.
	\begin{enumerate}
		\item If $G$ is decomposable, then $G \in$ \[ \{P_4 \circ D_{2,1}, D_{2,1} \circ P_4, P_4 \circ \overline{D_{2,1}}, \overline{D_{2,1}} \circ P_4, P_4^3, P_4 \circ C_4, P_4 \circ 2K_2\}, \]
		\item If $G$ is indecomposable, then
		\[ G \in \{T_{111}, \overline{T_{111}}, D_{3,1}, \overline{D_{3,1}}, T_{221}, \overline{T_{221}}\}, \]
		all of which are split (see  \Cref{grafos.primos.deg=3}).
	\end{enumerate}
\end{theorem}

%\begin{figure}[h]
%	\centering
%	\includegraphics[scale=0.8]{grafos_primos_de_grado_3.pdf}
%	\caption{The 6 prime graphs of degree $3$.}
%	\label{grafos.primos.deg=3}
%\end{figure}
\begin{proof}
	It follows from the previous discussion and \Cref{no.existe.deg3.primo}.
\end{proof}



%%%%%%%%%%%%%%%%%%%%%%%%%%%%%%%%%%%%%%%%%%%%%%%%%%%%%%%%%%%

\section{Prime split graphs of degree 4}\label{sec_deg=4}

We now classify the prime split graphs of degree 4. Consider a prime graph $G$ such that $\deg(G) = |\mathcal{Q}_G(P_4)| = 4$. Thanks to  \Cref{5>degG=|{P_4}|.implica.G.split}, we know that $G$ is split. Assume $G$ has bipartition $(K,I)$ and $\alpha = |I| \leq |K| = \omega$. Clearly, $\omega \geq 3$, since otherwise we would have $G\approx P_4$. Thus, $2 \leq \alpha \leq 5$ and $3 \leq \omega \leq 5$, by \Cref{|I|_|K|<=2deg(S)}. Applying \Cref{S.primo.iff.Phi(S).conexo}, we have that $\Phi = \Phi(G)$ is isomorphic to one of the 12 multigraphs in  \Cref{Phi.posibles.deg=4}.
\begin{figure}[h]
	\centering
	\begin{tikzpicture}[scale=1, every node/.style={draw, circle, fill=white, minimum size=0.25cm, inner sep=1pt}]
		
		% First row
		\node (a1) at (0,0) {};
		\node (a2) [right=of a1] {};
		\draw (a1) -- (a2);
		\draw[bend left=15] (a1) to (a2);
		\draw[bend left=35] (a1) to (a2);
		\draw[bend right=15] (a1) to (a2);
		
		\node (b1) [right=2cm of a2] {};
		\node (b2) [right=of b1] {};
		\node (b3) [above=of b1] {};
		\draw (b1) -- (b2);
		\draw (b1) -- (b3);
		\draw[bend left=20] (b1) to (b2);
		\draw[bend right=20] (b1) to (b3);
		
		\node (c1) [right=2cm of b2] {};
		\node (c2) [right=of c1] {};
		\node (c3) [above=of c1] {};
		\draw (c1) -- (c2);
		\draw (c1) -- (c3);
		\draw (c2) -- (c3);
		\draw[bend left=20] (c2) to (c3);
		
		% Second row
		\node (d1) at (0,-2) {};
		\node (d2) [right=of d1] {};
		\node (d3) [above=of d1] {};
		\node (d4) [above=of d2] {};
		\draw (d1) -- (d2);
		\draw (d1) -- (d3);
		\draw (d2) -- (d4);
		\draw (d3) -- (d4);
		
		\node (e1) [right=2cm of d2] {};
		\node (e2) [right=of e1] {};
		\node (e3) [above=of e1] {};
		\node (e4) [above=of e2] {};
		\draw (e4) -- (e3);
		\draw (e1) -- (e3);
		\draw[bend left=20] (e1) to (e3);
		\draw (e2) -- (e3);
		
		\node (f1) [right=2cm of e2] {};
		\node (f2) [right=of f1] {};
		\node (f3) [above=of f1] {};
		\draw (f1) -- (f3);
		\draw (f2) -- (f3);
		\draw[bend left=20] (f3) to (f1);
		\draw[bend right=20] (f3) to (f1);
		
		% Third row
		\node (g1) at (0,-4) {};
		\node (g2) [right=of g1] {};
		\node (g3) [above=of g1] {};
		\node (g4) [above=of g2] {};
		\draw (g1) -- (g2);
		\draw (g1) -- (g3);
		\draw (g3) -- (g4);
		\draw[bend left=20] (g3) to (g1);
		
		\node (h1) [right=2cm of g2] {};
		\node (h2) [right=of h1] {};
		\node (h3) [above=of h1] {};
		\node (h4) [above=of h2] {};
		\draw (h1) -- (h2);
		\draw (h1) -- (h3);
		\draw (h2) -- (h4);
		\draw (h2) -- (h3);
		
		\node (i1) [right=2cm of h2] {};
		\node (i2) [right=of i1] {};
		\node (i3) [above=of i1] {};
		\node (i4) [above=of i2] {};
		\draw (i1) -- (i2);
		\draw (i1) -- (i3);
		\draw (i2) -- (i4);
		\draw[bend left=20] (i1) to (i3);
		
		% Fourth row
		\node (j1) at (0,-6) {};
		\node (j2) [right=of j1] {};
		\node (j3) [above=of j1] {};
		\node (j4) [above=of j2] {};
		\node (j5) [below=of j1] {};
		\draw (j1) -- (j5);
		\draw (j1) -- (j3);
		\draw (j2) -- (j4);
		\draw (j3) -- (j4);
		
		\node (k1) [right=2cm of j2] {};
		\node (k2) [right=of k1] {};
		\node (k3) [above=of k1] {};
		\node (k4) [above=of k2] {};
		\node (k5) [below=of k1] {};
		\draw (k1) -- (k2);
		\draw (k1) -- (k3);
		\draw (k1) -- (k4);
		\draw (k1) -- (k5);
		
		\node (l1) [right=2cm of k2] {};
		\node (l2) [right=of l1] {};
		\node (l3) [above=of l1] {};
		\node (l4) [above=of l2] {};
		\node (l5) [below=of l1] {};
		\draw (l5) -- (l2);
		\draw (l1) -- (l3);
		\draw (l1) -- (l4);
		\draw (l1) -- (l5);
	\end{tikzpicture}
	\caption{The 12 connected multigraphs of size 4.}
	\label{Phi.posibles.deg=4}
\end{figure}
If $\alpha=2$, then $G\approx D_{4,1}$ or $D_{2,2}$ by \Cref{S_activo_|I|=2} (see \Cref{split.primos.deg=4.fig.1}). Thanks to \Cref{hechos.previos}, we can immediately rule out that $\Phi$ is the $K_3$ with a pendant vertex, $P_5$, or $\overline{D_{2,1}}$. 
%\begin{theorem}[\cite{schvollner2026induced}]
%	\label{prohibido.P_4.sigma_23=1}
%	If $S$ is a split graph, then $\Phi(S)$ cannot contain an induced path $v_1 v_2 v_3 v_4$ where $\sigma_{23}=1$.
%\end{theorem}
Applying Theorem 4.1 of \cite{schvollner2026induced}, we can also exclude that $\Phi = v_1v_2v_3v_4$ with $\sigma_{23} = 1$. If $\Phi = abca$ with $\sigma_{ac} = 2$, then by  \Cref{prop.basicas.sigma_uv} we would have $d_a = d_b = d_c = d_a \pm 1$, which is absurd.

If $\Phi \approx K_{4,1}$, let $v_5$ be the vertex of degree 4 in $K_{4,1}$ and $v_i$ its leaves, $i \in [4]$. Applying \Cref{hechos.previos}, we have that the $N_i=N_j$ for all $i,j\in [4]$, so $K = N_1 \cup N_5$. On the other hand, since $G$ has no universal vertices, we have $\varnothing = \bigcap_{j=1}^5 N_j = N_1 \cap N_5$, and consequently, $d_1 = 1 = d_5$, by  \Cref{prop.basicas.sigma_uv}. Then, $5 = \alpha \leq \omega = 2$, an absurdity.

If $\Phi = v_1v_2v_3v_4v_1$, then $N_1 = N_3$ and $N_2 = N_4$, by \Cref{hechos.previos}, so $K = N_1 \cup N_2$. Since $G$ has no universal vertices, we have $\varnothing = \bigcap_{i=1}^4 N_i = N_1 \cap N_2$, and consequently, $d_1 = 1 = d_2$, by  \Cref{prop.basicas.sigma_uv}. Then, $4 = \alpha \leq \omega = d_1 + d_2 = 2$, which is impossible.

If $\Phi = v_1v_2v_3v_4$ with $\sigma_{23} = 2$, then we can assume without loss of generality that $d_2 \leq d_3$. Then, $d_1 = d_2$, $d_3 = d_1 + 1 = d_4$, $N_2 \subseteq N_4$, and $N_1 \subseteq N_3, N_4$, by  \Cref{prop.basicas.sigma_uv}. Thus, $K = N_3 \cup N_4$. Since $G$ has no universal vertices, we have $\varnothing = \bigcap_{i=1}^4 N_i = N_1 \cap N_2$, and consequently, $d_1 = 1$, by  \Cref{prop.basicas.sigma_uv}. Thus, $d_3 = 2$ and $\eta_{34} = 1$, again by  \Cref{prop.basicas.sigma_uv}. Finally, we obtain $4 = \alpha \leq \omega = d_3 + d_4 - \eta_{34} = 3$, an absurdity.

If $E(\Phi) = \{2v_1v_2, v_2v_3, v_2v_4\}$ (a $K_{3,1}$ with a double-edge), then $d_3 = d_2 = d_4$, $|d_2 - d_1| = 1$, and $N_3 = N_4$, by  \Cref{prop.basicas.sigma_uv}. If $d_2 = d_1 + 1$, then $N_1 \subseteq N_3$, so $K = N_2 \cup N_3$. Since $G$ has no universal vertices, we have $\varnothing = \bigcap_{i=1}^4 N_i = N_1 \cap N_2$, and consequently, $d_1 = 1$, by  \Cref{prop.basicas.sigma_uv}. Thus, $\eta_{23} = 1$, again by  \Cref{prop.basicas.sigma_uv}. Finally, we obtain $4 = \alpha \leq \omega = d_2 + d_3 - \eta_{23} = 3$, which is impossible. The case $d_1 = d_2 + 1$ is reduced to an absurdity with similar arguments.


%If $I = \{a, b\}$, then $N_a \cap N_b = \varnothing$, since $S$ has no universal vertices. Thus, $\omega = d_a + d_b$ and $4 = \sigma_{ab} = d_a d_b$, by  \Cref{prop.basicas.sigma_uv}. With $d_a \leq d_b$, we conclude that $(d_a, d_b, \omega) \in \{(1, 4, 5), (2, 2, 4)\}$. The realizations of these triples are called $D_{4,1}$ and $D_{2,2}$, respectively; they are shown along with their complements in  \Cref{split.primos.deg=4.fig.1}.
\begin{figure}[ht]
	\centering
	\begin{tikzpicture}[scale=0.9, 
		every node/.style={draw, circle, fill=white, inner sep=0pt, minimum size=6.5pt},
		]
		
		%%%%%% D_{4,1} : clique K_5 + dos independientes %%%%%%
		\begin{scope}[shift={(0,3.5)}]
			% pentágono (clique K_5) centrado en (0,0), radio 1.3
			\node (t)  at (90:1.3)  {}; % top
			\node (ur) at (18:1.3)  {}; % upper-right
			\node (lr) at (-54:1.3) {}; % lower-right
			\node (ll) at (234:1.3) {}; % lower-left
			\node (ul) at (162:1.3) {}; % upper-left
			% independientes
			\node (v1) at (-2.4, 1.3)  {}; % arriba-izquierda
			\node (v2) at ( 2.2, -1.053) {}; % abajo-derecha
			% K_5: todas las aristas
			\draw (t)--(ur); \draw (ur)--(lr); \draw (lr)--(ll); \draw (ll)--(ul); \draw (ul)--(t);
			\draw (t)--(lr); \draw (t)--(ll); \draw (ur)--(ll); \draw (ur)--(ul); \draw (lr)--(ul);
			% v1 adyacente a t, ul, ur, ll  (no a lr)
			\draw (v1)--(t);
			\draw (v1)--(ul);
			\draw (v1)--(ur);
			\draw (v1)--(ll);
			% v2 adyacente a lr
			\draw (v2)--(lr);
			\node[draw=none] at (-2.0, -0.2) {$D_{4,1}$};
		\end{scope}
		
		%%%%%% D_{2,2} : clique K_4 + dos independientes %%%%%%
		\begin{scope}[shift={(5.5,3)}]
			\node (ul) at (-0.55, 0.55) {};
			\node (ur) at ( 0.55, 0.55) {};
			\node (lr) at ( 0.55,-0.55) {};
			\node (ll) at (-0.55,-0.55) {};
			\node (w1) at (-1.85, 0)    {};
			\node (w2) at ( 1.85, 0)    {};
			% K_4: todas las aristas
			\draw (ul)--(ur); \draw (ur)--(lr); \draw (lr)--(ll); \draw (ll)--(ul);
			\draw (ul)--(lr); \draw (ur)--(ll);
			% w1: a ul, ll
			\draw (w1)--(ul); \draw (w1)--(ll);
			% w2: a ur, lr
			\draw (w2)--(ur); \draw (w2)--(lr);
			\node[draw=none] at (0, 1.2) {$D_{2,2}$};
		\end{scope}
		
		%%%%%% \overline{D_{4,1}} %%%%%%
		\begin{scope}[shift={(0,0.4)}]
			% clique: {A, B}, indep: {p1, p2, p3, p4, p5}
			% A adyacente a p1,p2,p3,p4 y a B;  B adyacente a p5 y a A
			\node (A)  at (0, 0)      {}; % central
			\node (p1) at (-0.4, 1)    {}; % arriba
			\node (p2) at (-1.3, 0)   {}; % izquierda
			\node (p3) at ( 1.3, 0)   {}; % derecha
			\node (p4) at (-0.4,-0.9){}; % abajo-izq
			\node (B)  at ( 0.9,-0.9) {}; % segundo clique, abajo-der
			\node (p5) at ( 2.2,-0.9) {}; % pendiente de B
			\draw (A)--(p1);
			\draw (A)--(p2);
			\draw (A)--(p3);
			\draw (A)--(p4);
			\draw (A)--(B);
			\draw (B)--(p5);
			\node[draw=none] at (1, 0.8) {$\overline{D_{4,1}}$};
		\end{scope}
		
		%%%%%% \overline{D_{2,2}} %%%%%%
		\begin{scope}[shift={(5.5,0)}]
			% clique: {A, B}; indep: {q1, q2, q3, q4}
			\node (A)  at (0, 0.5)    {};
			\node (q1) at (-1.3, 0.5) {};
			\node (q2) at ( 1.3, 0.5) {};
			\node (B)  at ( 0.7,-0.5) {};
			\node (q3) at (-0.6,-0.5) {};
			\node (q4) at ( 2.0,-0.5) {};
			\draw (A)--(q1);
			\draw (A)--(q2);
			\draw (A)--(B);
			\draw (B)--(q3);
			\draw (B)--(q4);
			\node[draw=none] at (0, 1.2) {$\overline{D_{2,2}}$};
		\end{scope}
		
	\end{tikzpicture}
	\caption{The prime split graphs $D_{4,1}, D_{2,2}$ and their complements.}
	\label{split.primos.deg=4.fig.1}
\end{figure}
%\begin{figure}[h]
%	\centering
%	\includegraphics[scale=0.8]{grafos_split_primos_deg=4_fig1.pdf}
%	\caption{The prime split graphs $D_{4,1}, D_{2,2}$ and their complements.}
%	\label{split.primos.deg=4.fig.1}
%\end{figure}

If $\Phi=abc$ ignoring multiplicities, then, exchanging $a$ and $c$ if necessary, \Cref{S_activo_|I|=3} gives $\sigma_{ab}=\kappa_b(\pi_b+\kappa_a)$ and $\sigma_{bc}=\pi_b(\pi_c+\kappa_b)$, with $\sigma_{ab}+\sigma_{bc}=4$ and $\pi_b,\kappa_b\geq 1$. A direct computation yields exactly the five graphs of the following table, where only the parameters of \Cref{S_activo_|I|=3} not forced to vanish are listed; as in \Cref{sec_deg=3}, the subscript of each name records the degrees of the $I$-vertices. In all cases $\omega\geq 3=\alpha$.
\begin{center}
	\begin{tabular}{cccc}
		\hline
		$(\sigma_{ab},\sigma_{bc})$ & $(\pi_b,\pi_c,\kappa_a,\kappa_b)$ & $\omega$ & $G$ \\
		\hline
		$(1,3)$ & $(1,2,0,1)$ & $4$ & $R_{311}$ \\
		$(3,1)$ & $(1,0,2,1)$ & $4$ & $R_{331}$ \\
		$(2,2)$ & $(1,1,1,1)$ & $4$ & $R_{321}$ \\
		$(2,2)$ & $(2,0,0,1)$ & $3$ & $R_{211}$ \\
		$(2,2)$ & $(1,0,0,2)$ & $3$ & $\overline{R_{211}}$ \\
		\hline
	\end{tabular}
\end{center}
These graphs, together with $\overline{R_{311}}$, $\overline{R_{331}}$ and $\overline{R_{321}}$, are shown in \Cref{split.primos.deg=4.fig.2,split.primos.deg=4.fig.3}.
%If $\Phi = abc$ with $\sigma_{bc} = 3$, then $d_a = d_b$ and $|d_c - d_b| = 2$, by \Cref{prop.basicas.sigma_uv}. If $d_c = d_b + 2$, it follows by \Cref{prop.basicas.sigma_uv} that $N_a \subseteq N_c$, so $K = N_b \cup N_c$. Since $G$ has no universal vertices, we have $\varnothing = N_a \cap N_b \cap N_c = N_a \cap N_b$, and consequently, $d_a = 1$, by \Cref{prop.basicas.sigma_uv}. Thus, $\eta_{bc} = 0$, again by \Cref{prop.basicas.sigma_uv}. Then, $\omega = d_b + d_c = 4$, so we have a realization for $G$ called $R_{311}$ (the subscript records the degrees of the $I$-vertices). If $d_b = d_c + 2$, it follows by  \Cref{prop.basicas.sigma_uv} that $N_c \subseteq N_a$, so $K = N_a \cup N_b$. Since $G$ has no universal vertices, we have $\varnothing = N_a \cap N_b \cap N_c = N_b \cap N_c$, and consequently, $d_c = 1$, by  \Cref{prop.basicas.sigma_uv}. Thus, $\eta_{ab} = 2$, again by  \Cref{prop.basicas.sigma_uv}. Then, $\omega = d_a + d_b - \eta_{ab} = 4$, so we obtain another realization for $G$ called $R_{331}$. The graphs $R_{311}, R_{331}$ and their complements are shown in  \Cref{split.primos.deg=4.fig.2}.
\begin{figure}[h]
	\centering
	\begin{tikzpicture}[scale=0.9, 
		every node/.style={draw, circle, fill=white, inner sep=0pt, minimum size=6.5pt},
		]
		
		%%%%%% R_{311} : clique K_4 + tres independientes %%%%%%
		\begin{scope}[shift={(1,3.2)}]
			\node (ul) at (0, 1)  {};
			\node (ur) at (1, 1)  {};
			\node (ll) at (0, 0)  {};
			\node (lr) at (1, 0)  {};
			\node (w1) at (-1.3, 1)   {}; % conectado a ul
			\node (w2) at (-1.3, 0)   {}; % conectado a ll
			\node (w3) at ( 2, 2) {}; % conectado a ul, ur, lr
			% K_4: todas las aristas
			\draw (ul)--(ur); \draw (ur)--(lr); \draw (lr)--(ll); \draw (ll)--(ul);
			\draw (ul)--(lr); \draw (ur)--(ll);
			% pendants
			\draw (w1)--(ul);
			\draw (w2)--(ll);
			\draw (w3)--(ul); \draw (w3)--(ur); \draw (w3)--(lr);
			\node[draw=none] at (-1.0, 1.8) {$R_{311}$};
		\end{scope}
		
		%%%%%% R_{331} : clique K_4 + tres independientes %%%%%%
		\begin{scope}[shift={(6.6,3.2)}]
			\node (ul) at (0, 1)  {};
			\node (ur) at (1, 1)  {};
			\node (ll) at (0, 0)  {};
			\node (lr) at (1, 0)  {};
			\node (w1) at (-1, 2) {}; % conectado a ul, ur
			\node (w2) at (2, 2) {}; % conectado a ul, ur
			\node (w3) at (-1.3, 0)   {}; % conectado a ll
			% K_4: todas las aristas
			\draw (ul)--(ur); \draw (ur)--(lr); \draw (lr)--(ll); \draw (ll)--(ul);
			\draw (ul)--(lr); \draw (ur)--(ll);
			% pendants
			\draw (w1)--(ul); \draw (w1)--(ur);
			\draw (w2)--(ul); \draw (w2)--(ur);
			\draw (w3)--(ll); \draw (ll)--(w1); 
			\draw (lr)--(w2);
			\node[draw=none] at (-1.2, 0.8) {$R_{331}$};
		\end{scope}
		
		%%%%%% \overline{R_{311}} %%%%%%
		\begin{scope}[shift={(1,1)}]
			% triángulo K_3: A arriba, B abajo-izq, C abajo-der
			\node (A) at (0, 1)      {};
			\node (B) at (-0.55, 0)  {};
			\node (C) at ( 0.55, 0)  {};
			\node (p1) at (-1.3, 1)  {};
			\node (p2) at ( 1.3, 1)  {};
			\node (p3) at ( 1.7, 0)  {};
			\node (p4) at ( 0.55,-1.1) {};
			\draw (A)--(B); \draw (B)--(C); \draw (A)--(C);
			\draw (A)--(p1);
			\draw (A)--(p2);
			\draw (A)--(p3);   % cruza con C--p2
			\draw (B)--(p1);
			\draw (C)--(p2);   % cruza con A--p3
			\draw (C)--(p3);
			\draw (C)--(p4);
			\node[draw=none] at (-1.1,-0.8) {$\overline{R_{311}}$};
		\end{scope}
		
		%%%%%% \overline{R_{331}} %%%%%%
		\begin{scope}[shift={(6.6,1)}]
			\node (A) at (0, 1)      {};
			\node (B) at (-0.55, 0)  {};
			\node (C) at ( 0.55, 0)  {};
			\node (q1) at (-1.3, 1)  {};
			\node (q2) at ( 1.3, 1)  {};
			\node (q3) at ( 1.7, 0)  {};
			\node (q4) at ( 0.55,-1.1) {};
			\draw (A)--(B); \draw (B)--(C); \draw (A)--(C);
			\draw (A)--(q2);
			\draw (B)--(q1);
			\draw (C)--(q3);
			\draw (C)--(q4);
			\draw (C)--(q2);
			\node[draw=none] at (-1.1,-0.8) {$\overline{R_{331}}$};
		\end{scope}
		
	\end{tikzpicture}
	\caption{The prime split graphs $R_{311}, R_{331}$ and their complements.}
	\label{split.primos.deg=4.fig.2}
\end{figure}
%\begin{figure}[H]
%	\centering
%	\includegraphics[scale=0.8]{grafos_split_primos_deg=4_fig2.pdf}
%	\caption{The prime split graphs $R_{311}, R_{331}$ and their complements.}
%	\label{split.primos.deg=4.fig.2}
%\end{figure}

%Let $\Phi = abc$ with $\sigma_{ab} = 2=\sigma_{bc}$. Then $|d_b - d_a| = 1 = |d_c - d_b|$, by  \Cref{prop.basicas.sigma_uv}. If $d_b = d_a + 1$ and $d_c = d_a$, it follows by  \Cref{prop.basicas.sigma_uv} that $N_a = N_c$, so $K = N_a \cup N_b$. Since $G$ has no universal vertices, we have $\varnothing = N_a \cap N_b \cap N_c = N_a \cap N_b$, and consequently, $d_a = 1$, by  \Cref{prop.basicas.sigma_uv}. Then, $\omega = d_a + d_b = 3$, so we have a realization for $G$ called $R_{211}$. If $d_b = d_a + 1$ and $d_c = d_a + 2$, it follows by  \Cref{prop.basicas.sigma_uv} that $N_a \subseteq N_c$, so $K = N_b \cup N_c$. Since $G$ has no universal vertices, we have $\varnothing = N_a \cap N_b \cap N_c = N_a \cap N_b$, and consequently, $d_a = 1$, by  \Cref{prop.basicas.sigma_uv}. Thus, $\eta_{bc} = 1$, again by  \Cref{prop.basicas.sigma_uv}. Then, $\omega = d_b + d_c - \eta_{bc} = 4$, so we obtain another realization for $G$ called $R_{321}$. Finally, if $d_a = d_b + 1$ and $d_a = d_c$, then analogous reasoning yields $\overline{R_{211}}$. The graphs $R_{211}, R_{321}$ and their complements are shown in  \Cref{split.primos.deg=4.fig.3}.
\begin{figure}[h]
	\centering
	\begin{tikzpicture}[scale=0.9, 
		every node/.style={draw, circle, fill=white, inner sep=0pt, minimum size=6.5pt},
		]
		
		%%%%%% R_{211} : K_3 + tres independientes %%%%%%
		\begin{scope}[shift={(0,3.2)}]
			\node (T)  at (0, 1)      {};
			\node (BL) at (-0.55, 0)  {};
			\node (BR) at ( 0.55, 0)  {};
			\node (p1) at ( 1.4, 1.3) {}; % conectado a T
			\node (p2) at ( 1.9, 0.2) {}; % conectado a T (arista larga)
			\node (p3) at ( 0, -1.1)  {}; % conectado a BL, BR
			% K_3
			\draw (T)--(BL); \draw (T)--(BR); \draw (BL)--(BR);
			% externas
			\draw (T)--(p1);
			\draw (T)--(p2);
			\draw (BL)--(p3);
			\draw (BR)--(p3);
			\node[draw=none] at (1.3, -0.9) {$R_{211}$};
		\end{scope}
		
		%%%%%% R_{321} : K_4 + tres independientes %%%%%%
		\begin{scope}[shift={(6.6,2.2)}]
			\node (ul) at (0, 1) {};
			\node (ur) at (1, 1) {};
			\node (ll) at (0, 0) {};
			\node (lr) at (1, 0) {};
			\node (w1) at (-1.3, 1.8) {}; % conectado a ul, ur
			\node (w2) at ( 2.3, 1.8) {}; % conectado a ur, lr
			\node (w3) at (-1.3, 0)   {}; % conectado a ll
			% K_4
			\draw (ul)--(ur); \draw (ur)--(lr); \draw (lr)--(ll); \draw (ll)--(ul);
			\draw (ul)--(lr); \draw (ur)--(ll);
			% externas
			\draw (w1)--(ul); \draw (w1)--(ur);
			\draw (w2)--(ur); \draw (w2)--(lr);
			\draw (w3)--(ll);
			\draw (w1)--(ll);
			\node[draw=none] at (2.3, 0.1) {$R_{321}$};
		\end{scope}
		
		%%%%%% \overline{R_{211}} %%%%%%
		\begin{scope}[shift={(0,0)}]
			\node (A) at (0, 1)      {};
			\node (B) at (-0.55, 0)  {};
			\node (C) at ( 0.55, 0)  {};
			\node (q1) at ( 1.4, 1)  {};
			\node (q2) at ( 1.7, 0)  {};
			\node (q3) at (0,-0.9){};
			\draw (A)--(B); \draw (B)--(C); \draw (A)--(C);
			% se cruzan A--q2 y C--q1
			\draw (A)--(q1);
			\draw (A)--(q2);
			\draw (C)--(q1);
			\draw (C)--(q2);
			\draw (B)--(q3);
			\node[draw=none] at (1.4, -0.9) {$\overline{R_{211}}$};
		\end{scope}
		
		%%%%%% \overline{R_{321}} %%%%%%
		\begin{scope}[shift={(6.6,0)}]
			\node (A) at (0, 1)       {};
			\node (B) at (-0.55, 0)   {};
			\node (C) at ( 0.55, 0)   {};
			\node (q1) at (-1.3, 1)   {};
			\node (q2) at ( 1.3, 1)   {};
			\node (q3) at ( 1.7, 0)   {};
			\node (q4) at ( 0, -1.1)  {};
			\draw (A)--(B); \draw (B)--(C); \draw (A)--(C);
			\draw (A)--(q2);
			\draw (B)--(q1);
			\draw (C)--(q3);
			\draw (B)--(q4);
			\draw (C)--(q4);
			\draw (C)--(q2);
			\node[draw=none] at (2.4, -0.9) {$\overline{R_{321}}$};
		\end{scope}
		
	\end{tikzpicture}
	\caption{The prime split graphs $R_{211}, R_{321}$ and their complements.}
	\label{split.primos.deg=4.fig.3}
\end{figure}
%\begin{figure}[H]
%	\centering
%	\includegraphics[scale=0.8]{grafos_split_primos_deg=4_fig3.pdf}
%	\caption{The prime split graphs $R_{211}, R_{321}$ and their complements.}
%	\label{split.primos.deg=4.fig.3}
%\end{figure}

\begin{theorem}
	\label{clasificacion.split.primos.deg=4}
	The 12 graphs in \Cref{split.primos.deg=4.fig.1,split.primos.deg=4.fig.2,split.primos.deg=4.fig.3} are all the prime split graphs of degree 4.
\end{theorem}



\begin{proof}
	It follows from the previous discussion.
\end{proof}		

%%%%%%%%%%%%%%%%%%%%%%%%%%%%%%%%%%%%%%%%%%%%%%%%%%%%%%%%%%%%%%


\section{Prime non-split graphs of degree 4}\label{sec_deg=4.no.split}


For $j\in[3]$, let $B_j$ be the graph with vertex set $\{v,c_0,c_1,c_2,c_3\}$ in which $c_0c_1c_2c_3c_0$ is an induced $C_4$ and $N_{B_j}(v)=\{c_0,\dots,c_{j-1}\}$. 
%Thus $B_1$ is the banner, $B_2$ is the
%house, and $B_3$ is obtained from $K_{2,3}$ by adding an edge joining
%two vertices of its part of size three; moreover
%$\overline{B_2}\approx P_5$ and
%$\overline{B_3}\approx P_3\,\dot\cup\,K_2$. 
These six graphs are displayed in \Cref{no.split.primos.deg=4.fig}. In this section we show that they are all the prime non-split graphs of degree 4.

\begin{figure}[ht]
	\centering
	\begin{tikzpicture}[scale=0.9,
		every node/.style={draw, circle, fill=white, inner sep=0pt,
			minimum size=6.5pt},
		]
		
		%%%%%% B_1 : the banner %%%%%%
		\begin{scope}[shift={(0,0)}]
			\node (c0) at (0,1) {};
			\node (c1) at (1,1) {};
			\node (c2) at (1,0) {};
			\node (c3) at (0,0) {};
			\node (v)  at (-0.9,1.8) {};
			\draw (c0)--(c1); \draw (c1)--(c2); \draw (c2)--(c3); \draw (c3)--(c0);
			\draw (v)--(c0);
			\node[draw=none] at (0.5,-0.85) {$B_1$};
		\end{scope}
		
		%%%%%% B_2 : the house %%%%%%
		\begin{scope}[shift={(4,0)}]
			\node (c0) at (0,1) {};
			\node (c1) at (1,1) {};
			\node (c2) at (1,0) {};
			\node (c3) at (0,0) {};
			\node (v)  at (0.5,1.9) {};
			\draw (c0)--(c1); \draw (c1)--(c2); \draw (c2)--(c3); \draw (c3)--(c0);
			\draw (v)--(c0); \draw (v)--(c1);
			\node[draw=none] at (0.5,-0.85) {$B_2$};
		\end{scope}
		
		%%%%%% B_3 %%%%%%
		\begin{scope}[shift={(8,0)}]
			\node (c0) at (0,1) {};
			\node (c1) at (1,1) {};
			\node (c2) at (1,0) {};
			\node (c3) at (0,0) {};
			\node (v)  at (1.9,1.8) {};
			\draw (c0)--(c1); \draw (c1)--(c2); \draw (c2)--(c3); \draw (c3)--(c0);
			\draw (v)--(c0); \draw (v)--(c1); \draw (v)--(c2);
			\node[draw=none] at (0.5,-0.85) {$B_3$};
		\end{scope}
		
		%%%%%% complement of B_1 : triangle with a path of length 2 %%%%%%
		\begin{scope}[shift={(0,-3)}]
			\node (a) at (-0.6,1.0) {};
			\node (b) at (-0.6,0.0) {};
			\node (c) at (0.27,0.5) {};
			\node (d) at (1.14,0.5) {};
			\node (e) at (2.01,0.5) {};
			\draw (a)--(b); \draw (a)--(c); \draw (b)--(c);
			\draw (c)--(d); \draw (d)--(e);
			\node[draw=none] at (0.5,-0.6) {$\overline{B_1}$};
		\end{scope}
		
		%%%%%% complement of B_2 : P_5 %%%%%%
		\begin{scope}[shift={(4,-3)}]
			\node (p1) at (-0.9,0.5) {};
			\node (p2) at (-0.2,0.5) {};
			\node (p3) at ( 0.5,0.5) {};
			\node (p4) at ( 1.2,0.5) {};
			\node (p5) at ( 1.9,0.5) {};
			\draw (p1)--(p2); \draw (p2)--(p3); \draw (p3)--(p4); \draw (p4)--(p5);
			\node[draw=none] at (0.5,-0.6) {$\overline{B_2}$};
		\end{scope}
		
		%%%%%% complement of B_3 : P_3 + K_2 %%%%%%
		\begin{scope}[shift={(8,-3)}]
			\node (q1) at (-0.2,1.0) {};
			\node (q2) at ( 0.5,1.0) {};
			\node (q3) at ( 1.2,1.0) {};
			\node (r1) at ( 0.15,0.0) {};
			\node (r2) at ( 0.85,0.0) {};
			\draw (q1)--(q2); \draw (q2)--(q3); \draw (r1)--(r2);
			\node[draw=none] at (0.5,-0.6) {$\overline{B_3}$};
		\end{scope}
		
	\end{tikzpicture}
	\caption{The six non-split prime graphs of degree 4.}
	\label{no.split.primos.deg=4.fig}
\end{figure}


\begin{theorem}
	\label{clasificacion.no.split.primos.deg=4}
	The graphs $B_1,B_2,B_3$ and their complements (see \Cref{no.split.primos.deg=4.fig}) are all the non-split prime graphs of degree $4$.
\end{theorem}



\begin{proof}
		Fix $j\in[3]$ and write $C=B_j[c_0,c_1,c_2,c_3]$. As $|B_j|=5$, every
	member of $\mathcal{Q}_{B_j}$ either equals $C$ or contains $v$; hence
	$\mathcal{Q}_{B_j}=\{C\}\cup\mathcal{H}_v$, so $\deg(B_j)=2+2=4$ by
	\Cref{conteo.local.C4} and \eqref{deg.como.suma}, and $B_j$ is active.
	If $B_j=S\circ H$ with $V(S),V(H)\neq\varnothing$, then $S$ and $H$
	are active by \Cref{hechos.previos}, so $|S|,|H|\geq 4$ by
	\Cref{|G^*|} and $|B_j|\geq 8$, which is absurd; hence $B_j$ is prime.
	Moreover $B_j$ is not split, because $C\approx C_4$. Since
	complementation preserves the degree, primality and non-splitness,
	the same holds for $\overline{B_1},\overline{B_2},\overline{B_3}$.
	Finally, these six graphs are clearly pairwise non-isomorphic.
	
	Conversely, let $G$ be a prime graph with $\deg(G)=4$ which is not
	split. By \Cref{5>degG=|{P_4}|.implica.G.split} we have
	$\deg(G)\neq|\mathcal{Q}_G(P_4)|$, so $\mathcal{Q}_G(C_4)\cup \mathcal{Q}_G(2K_2)\neq\varnothing$ by
	\Cref{degreeofG}. Complementation preserves the degree, primality and
	non-splitness, and interchanges induced $C_4$'s with induced
	$2K_2$'s; as the six graphs of the statement form three complementary
	pairs, we may assume $\mathcal{Q}_G(C_4)\neq\varnothing$ and it suffices to
	prove that $G\approx B_j$ for some $j\in[3]$. Fix an induced $C=c_0c_1c_2c_3c_0\approx C_4$ in
	$G$ and put $M=V(C)$.
	
		By \Cref{C4.distinguidores}, at most one vertex $v_0\in V(G)-M$
	distinguishes $M$, and $N_{v_0}\cap M$ is not a pair of opposite
	vertices of $C$; by \Cref{conteo.local.C4},
	$\sum_{H\in\mathcal{H}_{v_0}}\deg(H)=2$.
	
	Suppose first that no such $v_0$ exists. Then no vertex of $V(G)-M$
	distinguishes $M$, so $M$ is a module of $G$; moreover $G\neq C$
	because $\deg(C)=2$. This contradicts \Cref{C4.modulo.deg>=6}. Hence $v_0$ exists, and \eqref{conteo.global.C4} becomes an equality,
	so $\mathcal{Q}_G=\{C\}\cup\mathcal{H}_{v_0}$. In particular every member
	of $\mathcal{Q}_G$ has all its vertices in $\{v_0\}\cup M$, and since $G$
	is active we get $V(G)=\{v_0\}\cup M$. Finally, after a rotation of
	$C$ we may assume $N_{v_0}\cap M=\{c_0,\dots,c_{j-1}\}$ for some
	$j\in[3]$, because $N_{v_0}\cap M$ is neither empty, nor $M$, nor a
	pair of opposite vertices of $C$. Therefore $G\approx B_j$.
\end{proof}

\begin{corollary}
	\label{clasificacion.primos.deg=4}
	There are exactly 18 prime graphs of degree $4$: the 12
	split graphs of
	\Cref{split.primos.deg=4.fig.1,split.primos.deg=4.fig.2,split.primos.deg=4.fig.3},
	together with the 6 non-split graphs of \Cref{no.split.primos.deg=4.fig}.
\end{corollary}

%%%%%%%%%%%%%%%%%%%%%%%%%%%%%%%%%%%%%%%%%%%%%%%%%%%%%%%%%%%%%


\section{Consequences for realization graphs}
\label{sec_realization}

We now apply the previous classifications to unigraphs and realization
graphs; if $X$ has degree sequence $d$, we also write $\mathcal{G}(X)$ for
$\mathcal{G}(d)$. Following \cite{tyshkevich2000decomposition}, a degree
sequence is \emph{unigraphical} if all its realizations are isomorphic,
and a \emph{unigraph} is a realization of a unigraphical sequence;
following \cite{barrus2012switches}, a graph is a \emph{hereditary
	unigraph} if all its induced subgraphs are unigraphs, and it is
\emph{$\mathcal{F}$-free} if none of its induced subgraphs lies in the
family $\mathcal{F}$. Let
\[
\mathcal{F}=\bigl\{H,\overline{H}:\ H\in\{P_5,\ K_2\,\dot\cup\,K_3,\ B_1,\ 2P_3,\ K_2\,\dot\cup\,P_4,\ K_2\,\dot\cup\,C_4,\ \overline{T_{221}},\ R_{211}\}\bigr\}.
\]
%whose generators have degrees $4,6,4,8,7,10,3$ and $4$, respectively.
% by
%\Cref{degreeofG,clasificacion.activo_deg=3,clasificacion.split.primos.deg=4}.
In \cite{barrus2012switches} $\overline{T_{221}}$ and $R_{211}$ are called $R$ and
$S$ (see Figure~6 of
\cite{barrus2012switches}).
%thus the four members of
%$\mathcal{F}$ without a standard name are exactly
%$T_{221},\overline{T_{221}},R_{211},\overline{R_{211}}$, which arise on
%their own in our classification 

\begin{theorem}[\cite{barrus2012switches}, Theorem 4.2]
	\label{barrus.hereditarios}
	A graph is a hereditary unigraph if and only if it is $\mathcal{F}$-free.
\end{theorem}

If $H\prec G$, then $\mathcal{Q}_H\subseteq\mathcal{Q}_G$, so
\eqref{deg.como.suma} gives
\begin{equation}
	\label{deg.monotono}
	H\prec G\ \Longrightarrow\ \deg(H)\leq\deg(G).
\end{equation}

\begin{theorem}
	\label{grado<=3.unigrafo}
	Let $\deg(G)\leq 3$. Then $G$ is a unigraph if and only if it is a
	hereditary unigraph, if and only if it contains neither $T_{221}$ nor
	$\overline{T_{221}}$ as an induced subgraph; and otherwise some
	2-switch $\tau$ on $G$ satisfies $\deg(\tau(G))\geq 4$. In particular,
	every graph of degree at most $2$ is a hereditary unigraph.
\end{theorem}

\begin{proof}
	As $\deg(\overline{H})=\deg(H)$ for every $H$ \cite{pastine20252}, the
	members of $\mathcal{F}$ other than $T_{221},\overline{T_{221}}$ have
	degree at least $4$. Hence, if $T_{221},\overline{T_{221}}\not\prec G$,
	then $G$ is $\mathcal{F}$-free by \eqref{deg.monotono}, so it is a
	hereditary unigraph by \Cref{barrus.hereditarios}, and in particular a
	unigraph; the same argument proves the last assertion, because
	\eqref{deg.monotono} forbids $T_{221},\overline{T_{221}}$ inside a
	graph of degree at most $2$.
	
	Assume now $T_{221}\prec G$ or $\overline{T_{221}}\prec G$. A 2-switch
	on $\overline{G}$ is a 2-switch on $G$ on the same 4 vertices, the
	resulting graphs being complementary, so we may assume
	$T_{221}\prec G$. It is easy to check that there is a 2-switch $\tau$ on $T_{221}$ such that $\tau(T_{221})\approx \overline{R_{211}}$. Since $\tau(T_{221})\prec\tau(G)$, we have 
%	By the construction of $T_{221}$ in
%	\Cref{sec_deg=3}, write $V(T_{221})=\{x_1,x_2,x_3,a,b,c\}$ with
%	$\{x_1,x_2,x_3\}$ a clique and, inside $T_{221}$, $N_a=\{x_1\}$,
%	$N_b=\{x_2,x_3\}$, $N_c=\{x_1,x_2\}$. Then $ax_1x_3b\prec G$; let
%	$\tau$ replace $ax_1,x_3b$ with $ax_3,x_1b$. 
%	As $\tau$ changes no pair
%	outside $V(T_{221})$, the graph $\tau(G)[V(T_{221})]$ is the split
%	graph with clique $\{x_1,x_2,x_3\}$ and $N_a=\{x_3\}$,
%	$N_b=N_c=\{x_1,x_2\}$, that is, $\overline{R_{211}}$ (see
%	\Cref{split.primos.deg=4.fig.3}). 
	$\deg(\tau(G))\geq 4$ by
	\eqref{deg.monotono}, and since $\tau(G)$ and $G$ share their degree
	sequence, $G$ is not a unigraph.
\end{proof}

\begin{corollary}
	\label{G(d).transitivo}
	Let $X$ realize $d$, with $\deg(X)=k\leq 3$. Then $\mathcal{G}(d)$ is regular if and
	only if $T_{221},\overline{T_{221}}\not\prec X$, and in that case $d$ is unigraphical,
	$\mathcal{G}(d)$ is vertex-transitive and $k$-regular, and $\tau(X)\approx X$ for every
	2-switch $\tau$ on $X$.
\end{corollary}

\begin{proof}
	By \Cref{grado<=3.unigrafo}, $d$ is unigraphical. Let $\Gamma$ be the
	group of permutations $\pi$ of $V(X)$ with $d_{\pi(v)}=d_v$ for all
	$v$. Relabeling by $\pi\in\Gamma$ sends realizations of $d$ to
	realizations of $d$ and preserves the relation of differing by a
	2-switch, so $\Gamma$ acts on $\mathcal{G}(d)$ by automorphisms; and
	any isomorphism between realizations of $d$ preserves degrees, hence
	lies in $\Gamma$. So $\Gamma$ is transitive, $\mathcal{G}(d)$ is
	vertex-transitive, and it is $k$-regular because $X$ has degree $k$ in
	it. The last assertion follows from \Cref{grado<=3.unigrafo}.
\end{proof}

Thus $\mathcal{G}(T_{221})\approx K_{1,2,2}$: four of its vertices are
isomorphic to $T_{221}$ and have degree $3$, while the fifth is
isomorphic to $\overline{R_{211}}$ and has degree $4$; so $T_{221}$ and
$\overline{R_{211}}$ are the two realizations of $d(T_{221})$ up to
isomorphism. This is the first instance of a phenomenon occurring in
every degree $k\geq 3$.

\begin{proposition}
	\label{familia.Gk}
	For every $k\geq 3$ there are a prime graph $G_k$ with $\deg(G_k)=k$
	and a 2-switch $\tau$ on $G_k$ with $\deg(\tau(G_k))=2k-2$; in
	particular, $\mathcal{G}(G_k)$ is not regular.
\end{proposition}

\begin{proof}
	Let $G_k$ be the split graph with clique $K=\{x,y_1,\dots,y_{k-1}\}$,
	independent set $I=\{a,b,c\}$, and $N_a=\{x\}$,
	$N_b=\{y_1,\dots,y_{k-1}\}$, $N_c=\{y_1\}$. By
	\Cref{prop.basicas.sigma_uv}, $\sigma_{ab}=k-1$, $\sigma_{ac}=1$ and
	$\sigma_{bc}=0$, so $\deg(G_k)=k$ and $\Phi(G_k)=cab$ ignoring
	multiplicities. Every vertex of $K\cup\{a,b\}$ lies in an induced $P_4$
	$axy_ib$, and $c$ in the induced $P_4$ $axy_1c$; hence $G_k$ is active,
	and prime by \Cref{S.primo.iff.Phi(S).conexo}. Let $\tau$ replace
	$ax,by_1$ with $ay_1,bx$. In $\tau(G_k)$ we have $N_a=N_c=\{y_1\}$ and
	$N_b=\{x,y_2,\dots,y_{k-1}\}$, so $\sigma_{ab}=k-1=\sigma_{cb}$ and
	$\sigma_{ac}=0$, whence $\deg(\tau(G_k))=2k-2\neq k$.
\end{proof}

Note that $G_3\approx\overline{T_{221}}$ and $G_4\approx R_{311}$ (see
\Cref{grafos.primos.deg=3,split.primos.deg=4.fig.2}). In \Cref{R.k-1.1.1}, $G_5$ and $\tau(G_5)$ are shown.

\begin{figure}[h]
	\centering
	\begin{tikzpicture}[scale=0.8,
			every node/.style={draw, circle, fill=white, inner sep=0pt, minimum size=6.5pt},
			lbl/.style={draw=none, fill=none},
			sw/.style={thick},
			]
			
			%%%%%% $R_{4,1,1}$ %%%%%%
			\begin{scope}[shift={(0,0)}]
					\node (A1)[fill=gray] at (0,2.7) {};
					\node (A2)[fill=gray] at (150:1.4) {};
					\node (A3) at (110:1.4) {};
					\node (A4) at (70:1.4) {};
					\node (A5) at (30:1.4) {};
					\foreach \i in {1,...,5}
					\foreach \j in {1,...,5}
					{\ifnum\i<\j \draw (A\i)--(A\j);\fi}
					\node (Aa)[fill=gray] at (-0.85,2.7) {};
					\node (Ab)[fill=gray] at (0,0) {};
					\node (Ac) at (-2.5,1.15) {};
					\draw (Ac)--(A2);
					\draw[sw] (Aa)--(A1);
					\draw[sw] (Ab)--(A2);
					\draw (Ab)--(A3); \draw (Ab)--(A4); \draw (Ab)--(A5);
					\node[lbl] at (-1.3,2.7) {$a$};
					\node[lbl] at (0,-0.5) {$b$};
					\node[lbl] at (-2.98,1.15) {$c$};
					\node[lbl] at (0.38,2.78) {$x$};
					\node[lbl] at (-1.5,0.4) {$y_1$};
					\node[lbl] at (-1.5,-0.5) {$G_5$};
				\end{scope}
			
			%%%%%% $\tau(R_{4,1,1})$ %%%%%%
			\begin{scope}[shift={(6.6,0)}]
					\node (B1)[fill=gray] at (0,2.7) {};
					\node (B2)[fill=gray] at (150:1.4) {};
					\node (B3) at (110:1.4) {};
					\node (B4) at (70:1.4) {};
					\node (B5) at (30:1.4) {};
					\foreach \i in {1,...,5}
					\foreach \j in {1,...,5}
					{\ifnum\i<\j \draw (B\i)--(B\j);\fi}
					\node (Ba)[fill=gray] at (-0.85,2.7) {};
					\node (Bb)[fill=gray] at (0,0) {};
					\node (Bc) at (-2.5,1.15) {};
					\draw (Bc)--(B2);
					\draw[sw] (Ba)--(B2);
					\draw[sw] (Bb) to (B1);
					\draw (Bb)--(B3); \draw (Bb)--(B4); \draw (Bb)--(B5);
					\node[lbl] at (-1.3,2.7) {$a$};
					\node[lbl] at (0,-0.5) {$b$};
					\node[lbl] at (-2.98,1.15) {$c$};
					\node[lbl] at (0.38,2.78) {$x$};
					\node[lbl] at (-1.5,0.4) {$y_1$};
					\node[lbl] at (1.5,-0.5) {$\tau(G_5)$};
				\end{scope}
			
			\draw[-{Latex[length=2.2mm]}, thick] (1.85,1.35) -- (3.25,1.35);
			\node[lbl] at (2.55,1.72) {$\tau$};
			
	%		%%%%%% $R_{5,1,1}$ %%%%%%
	%		\begin{scope}[shift={(0,-6.3)}]
		%			\node (C1) at (0,2.7) {};
		%			\node (C2) at (150:1.4) {};
		%			\node (C3) at (120:1.4) {};
		%			\node (C4) at (90:1.4) {};
		%			\node (C5) at (60:1.4) {};
		%			\node (C6) at (30:1.4) {};
		%			\foreach \i in {1,...,6}
		%			\foreach \j in {1,...,6}
		%			{\ifnum\i<\j
			%				\ifnum\i=1
			%				\ifnum\j=2\else\ifnum\j=6\else\draw (C\i)--(C\j);\fi\fi
			%				\else\draw (C\i)--(C\j);\fi
			%				\fi}
		%			\draw (C1) to[bend right=13] (C2);
		%			\draw (C1) to[bend left=13] (C6);
		%			\node (Ca) at (-0.85,2.7) {};
		%			\node (Cb) at (0,0) {};
		%			\node (Cc) at (-2.5,1.15) {};
		%			\draw (Cc)--(C2);
		%			\draw[sw] (Ca)--(C1);
		%			\draw[sw] (Cb)--(C2);
		%			\draw (Cb)--(C3); \draw (Cb)--(C4); \draw (Cb)--(C5); \draw (Cb)--(C6);
		%			\node[lbl] at (-1.3,2.7) {$a$};
		%			\node[lbl] at (0,-0.5) {$b$};
		%			\node[lbl] at (-2.98,1.15) {$c$};
		%			\node[lbl] at (0.38,2.78) {\scriptsize $x$};
		%			\node[lbl] at (-1.5,0.4) {\scriptsize $y_1$};
		%			\node[lbl] at (0,-1.35) {$R_{5,1,1}$};
		%		\end{scope}
	%		
	%		%%%%%% $\tau(R_{5,1,1})$ %%%%%%
	%		\begin{scope}[shift={(6.6,-6.3)}]
		%			\node (D1) at (0,2.7) {};
		%			\node (D2) at (150:1.4) {};
		%			\node (D3) at (120:1.4) {};
		%			\node (D4) at (90:1.4) {};
		%			\node (D5) at (60:1.4) {};
		%			\node (D6) at (30:1.4) {};
		%			\foreach \i in {1,...,6}
		%			\foreach \j in {1,...,6}
		%			{\ifnum\i<\j
			%				\ifnum\i=1
			%				\ifnum\j=2\else\ifnum\j=6\else\draw (D\i)--(D\j);\fi\fi
			%				\else\draw (D\i)--(D\j);\fi
			%				\fi}
		%			\draw (D1) to[bend right=13] (D2);
		%			\draw (D1) to[bend left=13] (D6);
		%			\node (Da) at (-0.85,2.7) {};
		%			\node (Db) at (0,0) {};
		%			\node (Dc) at (-2.5,1.15) {};
		%			\draw (Dc)--(D2);
		%			\draw[sw] (Da)--(D2);
		%			\draw[sw] (Db) to[bend right=85] (D1);
		%			\draw (Db)--(D3); \draw (Db)--(D4); \draw (Db)--(D5); \draw (Db)--(D6);
		%			\node[lbl] at (-1.3,2.7) {$a$};
		%			\node[lbl] at (0,-0.5) {$b$};
		%			\node[lbl] at (-2.98,1.15) {$c$};
		%			\node[lbl] at (0.38,2.78) {\scriptsize $x$};
		%			\node[lbl] at (-1.5,0.4) {\scriptsize $y_1$};
		%			\node[lbl] at (0,-1.35) {$\tau(R_{5,1,1})$};
		%		\end{scope}
	%		
	%		\draw[-{Latex[length=2.2mm]}, thick] (1.85,-4.95) -- (3.25,-4.95);
	%		\node[lbl] at (2.55,-4.58) {$\tau$};
			
		\end{tikzpicture}
	\caption{The prime split graph $G_5$ and
			its image under the 2-switch $\tau$, which replaces $ax,by_1$ with
			$ay_1,bx$.}
	\label{R.k-1.1.1}
\end{figure}

We now turn to the realization graphs themselves. The first result below
reduces their study to prime graphs; in \cite{barrus2016realization} it is
stated for the canonical decomposition of $d$, whose components are the
degree sequences of $G_n,\dots,G_1$.

\begin{theorem}[\cite{barrus2016realization}, Theorem 6]
	\label{G.producto}
	If $G_n\circ\cdots\circ G_1$ is the Tyshkevich decomposition of $X$,
	then
	$\mathcal{G}(X)\approx\mathcal{G}(G_n)\,\square\cdots\square\,\mathcal{G}(G_1)$.
\end{theorem}

\begin{theorem}[\cite{barrus2023cliques}, page 2]
	\label{dual.spaces.iso}
	$\mathcal{G}(d)\approx\mathcal{G}(\overline{d})$.
\end{theorem}

The classification obtained in the previous sections determines the
realization graph of every prime graph of degree at most $4$; by
\Cref{dual.spaces.iso}, only one sequence in each complementary pair has
to be treated.

\begin{theorem}
	\label{realization.graphs.deg<5}
	Let $X$ be a prime graph with $\deg(X)\leq 4$. Then $\mathcal{G}(X)$
	is isomorphic to one of the ten graphs of
	\Cref{realization.graphs.fig}, according to the following table.
	\begin{center}
		\begin{tabular}{ll}
			\hline
			$\mathcal{G}(X)$ & $X$ \\
			\hline
			$K_2$ & $P_4$\\
			$K_3$ & $C_4$, $2K_2$, $D_{2,1}$, $\overline{D_{2,1}}$\\
			$K_4$ & $D_{3,1}$, $\overline{D_{3,1}}$\\
			$K_{3,3}$ & $T_{111}$, $\overline{T_{111}}$\\
			$K_{1,2,2}$ & $T_{221}$, $\overline{T_{221}}$, $R_{211}$, $\overline{R_{211}}$\\
			$K_5$ & $D_{4,1}$, $\overline{D_{4,1}}$\\
			$K_{2,2,2}$ & $D_{2,2}$, $\overline{D_{2,2}}$, $B_1$, $\overline{B_1}$, $B_3$, $\overline{B_3}$\\
			$K_{1,3,3}$ & $B_2$, $\overline{B_2}$\\
			$\mathcal{G}_1$ & $R_{311}$, $\overline{R_{311}}$, $R_{331}$, $\overline{R_{331}}$\\
			$\mathcal{G}_2$ & $R_{321}$, $\overline{R_{321}}$\\
			\hline
		\end{tabular}
	\end{center}
\end{theorem}

\begin{proof}
	By
	\Cref{clasificacion.activo_deg=2,clasificacion.activo_deg=3,clasificacion.split.primos.deg=4,clasificacion.no.split.primos.deg=4}
	there are exactly $29$ prime graphs of degree at most $4$, and they
	realize $26$ distinct degree sequences: the only coincidences are
	$d(T_{221})=d(\overline{R_{211}})$, $d(\overline{T_{221}})=d(R_{211})$
	and $d(B_1)=d(\overline{B_1})$. For each of these sequences, listing
	the realizations of $d$ and the 2-switches between them is a finite
	computation, which yields the table.
\end{proof}

\begin{figure}[ht]
	\centering
	\begin{tikzpicture}[scale=0.85,
		every node/.style={draw, circle, fill=white, inner sep=0pt, minimum size=6.5pt},
		nm/.style={draw=none, fill=none},
		]
		
		%%%%%%%%%%%%%%%%%%%% row 1 %%%%%%%%%%%%%%%%%%%%
		\begin{scope}[shift={(0,0)}]
			\node (a) at (0.15,0.9) {};
			\node (b) at (1.65,0.9) {};
			\draw (a)--(b);
			\node[nm] at (0.9,-0.45) {$K_2$};
		\end{scope}
		
		\begin{scope}[shift={(3.4,0)}]
			\node (a) at (0.9,1.7) {};
			\node (b) at (0.15,0.45) {};
			\node (c) at (1.65,0.45) {};
			\draw (a)--(b)--(c)--(a);
			\node[nm] at (0.9,-0.45) {$K_3$};
		\end{scope}
		
		\begin{scope}[shift={(6.8,0)}]
			\node (a) at (0.2,1.65) {};
			\node (b) at (1.6,1.65) {};
			\node (c) at (1.6,0.25) {};
			\node (d) at (0.2,0.25) {};
			\draw (a)--(b)--(c)--(d)--(a); \draw (a)--(c); \draw (b)--(d);
			\node[nm] at (0.9,-0.45) {$K_4$};
		\end{scope}
		
		\begin{scope}[shift={(10.2,0)}]
			\node (a) at (0.90,1.78) {};
			\node (b) at (0.09,1.19) {};
			\node (c) at (0.40,0.24) {};
			\node (d) at (1.40,0.24) {};
			\node (e) at (1.71,1.19) {};
			\draw (a)--(b)--(c)--(d)--(e)--(a);
			\draw (a)--(c); \draw (a)--(d); \draw (b)--(d); \draw (b)--(e); \draw (c)--(e);
			\node[nm] at (0.9,-0.45) {$K_5$};
		\end{scope}
		
		%%%%%%%%%%%%%%%%%%%% row 2 %%%%%%%%%%%%%%%%%%%%
		\begin{scope}[shift={(0,-3.3)}]
			\node (t1) at (0.15,1.65) {};
			\node (t2) at (0.90,1.65) {};
			\node (t3) at (1.65,1.65) {};
			\node (b1) at (0.15,0.25) {};
			\node (b2) at (0.90,0.25) {};
			\node (b3) at (1.65,0.25) {};
			\draw (t1)--(b1); \draw (t1)--(b2); \draw (t1)--(b3);
			\draw (t2)--(b1); \draw (t2)--(b2); \draw (t2)--(b3);
			\draw (t3)--(b1); \draw (t3)--(b2); \draw (t3)--(b3);
			\node[nm] at (0.9,-0.45) {$K_{3,3}$};
		\end{scope}
		
		\begin{scope}[shift={(3.4,-3.3)}]
			\node (a) at (0.15,1.65) {};
			\node (b) at (1.65,1.65) {};
			\node (c) at (1.65,0.25) {};
			\node (d) at (0.15,0.25) {};
			\node (o) at (0.90,0.95) {};
			\draw (a)--(b)--(c)--(d)--(a);
			\draw (o)--(a); \draw (o)--(b); \draw (o)--(c); \draw (o)--(d);
			\node[nm] at (0.9,-0.45) {$K_{1,2,2}$};
		\end{scope}
		
		\begin{scope}[shift={(6.8,-3.3)}]
			\node (h1) at (0.90,1.80) {};
			\node (h2) at (0.12,1.35) {};
			\node (h3) at (0.12,0.45) {};
			\node (h4) at (0.90,0.00) {};
			\node (h5) at (1.68,0.45) {};
			\node (h6) at (1.68,1.35) {};
			\draw (h1)--(h2)--(h3)--(h4)--(h5)--(h6)--(h1);
			\draw (h1)--(h3); \draw (h2)--(h4); \draw (h3)--(h5);
			\draw (h4)--(h6); \draw (h5)--(h1); \draw (h6)--(h2);
			\node[nm] at (0.9,-0.45) {$K_{2,2,2}$};
		\end{scope}
		
		\begin{scope}[shift={(10.2,-3.3)}]
			\node (h1) at (0.90,1.80) {};
			\node (h2) at (0.12,1.35) {};
			\node (h3) at (0.12,0.45) {};
			\node (h4) at (0.90,0.00) {};
			\node (h5) at (1.68,0.45) {};
			\node (h6) at (1.68,1.35) {};
			\node (o)  at (0.90,0.90) {};
			\draw (h1)--(h2)--(h3)--(h4)--(h5)--(h6)--(h1);
			\draw (o)--(h1); \draw (o)--(h2); \draw (o)--(h3);
			\draw (o)--(h4); \draw (o)--(h5); \draw (o)--(h6);
			\draw (h1) to[bend left=55] (h4);
			\draw (h2) to[bend right=55] (h5);
			\draw (h3) to[bend left=55] (h6);
			\node[nm] at (0.9,-0.45) {$K_{1,3,3}$};
		\end{scope}
		
		%%%%%%%%%%%%%%%%%%%% row 3 %%%%%%%%%%%%%%%%%%%%
		\begin{scope}[shift={(3.4,-6.6)}]
			\node (h1) at (0.90,1.80) {};
			\node (h2) at (0.12,1.35) {};
			\node (h3) at (0.12,0.45) {};
			\node (h4) at (0.90,0.00) {};
			\node (h5) at (1.68,0.45) {};
			\node (h6) at (1.68,1.35) {};
			\node (o)  at (0.90,0.90) {};
			\draw (h1)--(h2)--(h3)--(h4)--(h5)--(h6)--(h1);
			\draw (o)--(h1); \draw (o)--(h2); \draw (o)--(h3);
			\draw (o)--(h4); \draw (o)--(h5); \draw (o)--(h6);
			\draw (h1)--(h3); \draw (h4)--(h6);
			\draw (h2) to[bend right=55] (h5);
			\node[nm] at (0.9,-0.45) {$\mathcal{G}_1$};
		\end{scope}
		
		\begin{scope}[shift={(7.4,-6.6)}]
			\node (c1) at (0.00,1.90) {};
			\node (c2) at (1.90,1.90) {};
			\node (c3) at (1.90,0.00) {};
			\node (c4) at (0.00,0.00) {};
			\node (k1) at (0.95,1.45) {};
			\node (k2) at (1.60,0.95) {};
			\node (k3) at (0.95,0.45) {};
			\node (k4) at (0.30,0.95) {};
			\draw (c1)--(c2)--(c3)--(c4)--(c1);
			\draw (k1)--(k2)--(k3)--(k4)--(k1); \draw (k1)--(k3); \draw (k2)--(k4);
			\draw (c1)--(k1); \draw (c2)--(k1); \draw (c2)--(k2); \draw (c3)--(k2);
			\draw (c3)--(k3); \draw (c4)--(k3); \draw (c4)--(k4); \draw (c1)--(k4);
			\node[nm] at (0.95,-0.45) {$\mathcal{G}_2$};
		\end{scope}
		
	\end{tikzpicture}
	\caption{The ten realization graphs of \Cref{realization.graphs.deg<5}.}
	\label{realization.graphs.fig}
\end{figure}

The four graphs whose realization graph is $K_{1,2,2}$ are precisely the
four members of $\mathcal{F}$ drawn in Fig.~6 of
\cite{barrus2012switches}, and $T_{111}$ is the $3$-net, so
$\mathcal{G}(T_{111})\approx K_{3,3}$ is the transposition graph $T_3$, in
accordance with Example~3 of \cite{barrus2016realization}. In degree $4$
the converse of \Cref{G(d).transitivo} fails: $\mathcal{G}(B_1)\approx
K_{2,2,2}$ is regular, yet $d(B_1)=d(\overline{B_1})$ with
$B_1\not\approx\overline{B_1}$.

\begin{corollary}
	\label{G(d).grado.menor.5}
	If $\mathcal{G}(d)$ has a vertex $X$ of degree $k\in[4]$, then
	$\mathcal{G}(d)$ is the Cartesian product of the realization graphs of
	the prime factors of $X$, each of which is one of the ten graphs of
	\Cref{realization.graphs.fig}. In particular $\mathcal{G}(d)\approx
	K_2$ if $k=1$, and $\mathcal{G}(d)\approx K_3$ or $C_4$ if $k=2$.
\end{corollary}

\begin{proof}
	Let $G_n\circ\cdots\circ G_1$ be the Tyshkevich decomposition of $X$.
	If $|G_i|=1$, then $\mathcal{G}(G_i)\approx K_1$, a neutral factor for
	$\square$; otherwise $A_4(G_i)$ is connected by
	\Cref{indecomp.characterization}, so every vertex of $G_i$ lies in a
	member of $\mathcal{Q}_{G_i}$ and $G_i$ is prime. The prime factors
	$X_1,\dots,X_m$ of $X$ thus satisfy $\sum_i\deg(X_i)=k\leq 4$ by
	\Cref{hechos.previos}, and the first assertion follows from
	\Cref{G.producto,realization.graphs.deg<5}. If $k=1$, then $m=1$ and
	$X_1\approx P_4$; if $k=2$, then either $m=1$ and
	$\mathcal{G}(X_1)\approx K_3$, or $m=2$ and
	$\mathcal{G}(d)\approx K_2\,\square\,K_2\approx C_4$.
\end{proof}

%%%%%%%%%%%%%%%%%%%%%%%%%%%%%%%%%%%%%%%%%%%%%%%%%%%%%%%%%%%%%%%%%%%%%%%%%%%%

%%%%%%%%%%%%%%%%%%%%%%%%%%%%%%%%%%%%%%%%%%%%%%%%%%%%%%%%%%%%%%%%%%%%%%%%%%%%

%\section{Consequences for realization graphs}
%\label{sec_realization}
%
%In this section we collect the consequences of the previous classifications for realization graphs. The \emph{active part} $G^*$ of a graph $G$ is the subgraph of $G$ induced by its active vertices. 
%%We denote by $\act(G)$ the set of active vertices of $G$ and by $G^{*}=G[\act(G)]$ the \emph{active part} of $G$. 
%Since every member of $\mathcal{Q}_G$ is an induced subgraph of $G^{*}$, we have $\mathcal{Q}_{G^{*}}=\mathcal{Q}_G$; hence $G^{*}$ is active and $\deg(G^{*})=\deg(G)$. If $X$ is a graph with degree sequence $d$, we also write $\mathcal{G}(X)$ for $\mathcal{G}(d)$.
%
%\begin{theorem}[\cite{pastine20252}]
%	\label{isomorfismo.espacio.activo}
%	For every graph $X$, $\mathcal{G}(X)\approx\mathcal{G}(X^{*})$.
%\end{theorem}
%
%\begin{theorem}
%	\label{tau.preserves.deg1}
%	If $\deg(G)=1$, then $\deg(\tau(G))=1$ for the unique 2-switch $\tau$ on $G$. Moreover, $\tau(G)\approx G$.
%\end{theorem}
%
%\begin{proof}
%	By \Cref{|G^*|}, $G^{*}\approx P_4$, say $G^{*}=abcd$. Thus $\mathcal{G}(d)\approx\mathcal{G}(P_4)\approx K_2$ by \Cref{isomorfismo.espacio.activo}, and $\deg(\tau(G))=1$. It remains to show $\tau(G)\approx G$.
%	
%	Let $w\in V(G)-\{a,b,c,d\}$. Since $G^*$ is the only member of $\mathcal{Q}_G$, we have $G[w,a,b,c]\notin \mathcal{Q}_G$.	Recall $G[a,b,c]=abc$. If $wc\in E(G)$ and $wb\notin E(G)$, then $G[w,a,b,c]$ contains the edges $ab,bc,wc$ and lacks $ac,wb$; the remaining pair $wa$ then determines whether $G[w,a,b,c]$ is a $P_4$ (if $wa\notin E(G)$) or a $C_4$ (if $wa\in E(G)$). Both belong to $\mathcal{Q}_G$, contradicting the choice of $w$. Hence $wc\in E(G)\Rightarrow wb\in E(G)$. The same argument applied to $G[w,b,c,d]$ (which contains the path $bcd$) gives $wb\in E(G)\Rightarrow wc\in E(G)$. 
%	
%	Therefore, every $w\in V(G)-\{a,b,c,d\}$ is adjacent to both or to neither of $b,c$; equivalently, the transposition $\varphi:V(G)\to V(\tau(G))$ that swaps $b$ and $c$ and fixes every other vertex preserves adjacency (see Lemma 2.1 of \cite{barrus2012switches}), so $\tau(G)\approx G$. Since $\mathcal{G}(d)\approx K_2$, its two vertices are the only realizations of $d$, so $d$ has a unique realization up to isomorphism.
%\end{proof}
%
%\begin{corollary}
%	If $\mathcal{G}(d)$ has a leaf, then $\mathcal{G}(d) \approx K_2$.
%\end{corollary}
%
%\begin{proof}
%	It immediately follows from \Cref{tau.preserves.deg1}.
%\end{proof}
%
%\begin{theorem}
%	\label{2switch.preserva.deg=2}
%	Let $X$ be a graph. If $\deg(X)=2$ and $\tau$ is a 2-switch on $X$, then $\deg(\tau(X))=2$. In other words, the 2-switch preserves degree 2.
%\end{theorem}
%
%\begin{proof}
%	If $\deg(X)=2$, we know that $\deg(X^*)=2$ as well. By  \Cref{clasificacion.activo_deg=2}, $X^* \approx D_{2,1}$, $\overline{D_{2,1}}$, $C_4$, $2K_2$, or $P_4^2$ (see  \Cref{grafos.activos.deg<3}), and so it is straightforward to check that $\tau(X^*) \approx X^*$. Thus, $\mathcal{G}(X^*)$ is 2-regular. Finally, the result follows from  \Cref{isomorfismo.espacio.activo}.
%\end{proof}
%
%%\begin{theorem}[\cite{chartrand2010graphs}, page 24]
%%	\label{berge's.theorem}
%%	If $G$ and $H$ are two graphs with the same degree sequence, then there exists a sequence of 2-switches transforming $G$ into $H$.
%%\end{theorem}
%
%\begin{theorem}[\cite{barrus2023cliques}, page 2]
%	\label{dual.spaces.iso}
%	$\mathcal{G}(d)\approx\mathcal{G}(\overline{d})$.
%\end{theorem}
%
%\begin{corollary}
%	If $\mathcal{G}(d)$ has a vertex of degree 2, then $\mathcal{G}(d)\approx K_3$ or $C_4$.
%\end{corollary}
%
%\begin{proof}
%	From \Cref{2switch.preserva.deg=2} it follows that $\mathcal{G}(d)$ is 2-regular. Since $\mathcal{G}(d)$ is connected, we have that $\mathcal{G}(d)$ is a cycle. To complete the proof, thanks to \Cref{isomorfismo.espacio.activo,dual.spaces.iso}, it suffices to consider $X \in \{ \overline{D_{2,1}}, 2K_2, P_4^2 \}$ (see  \Cref{grafos.activos.deg<3}) and compute $c=|\mathcal{G}(X)|$. 
%		
%	If $X=\overline{D_{2,1}}$, we see that there are only 3 ways to label $X$ while maintaining its degree sequence, since the 2-switches on $X$ only exchange the adjacencies of its leaves. Thus, $c=3$. If $Q$ is a set of 4 elements, there are exactly 3 ways to write $Q$ as $A \dot{\cup} B$, with $|A|=2=|B|$. This shows that $c=3$ if $X=2K_2$. Finally, let $X=P_4^2$. The realization graph of a decomposable graph is the Cartesian product of the realization graphs of its Tyshkevich factors (see \cite{barrus2016realization}), so
%	\[
%	\mathcal{G}(P_4^2) \approx \mathcal{G}(P_4)\,\square\,\mathcal{G}(P_4) \approx K_2\,\square\,K_2 \approx C_4,
%	\]
%	and hence $c=4$.
%\end{proof}
%
%\begin{theorem}
%	\label{2switch.casi-preserva.deg3}
%	Let $X$ be a graph of degree 3 such that neither $T_{221}$ nor $\overline{T_{221}}$ (see \Cref{grafos.primos.deg=3}) are induced subgraphs of $X$. If $\tau$ is a 2-switch on $X$, then $\deg(\tau(X))=3$. In other words, $\mathcal{G}(X)$ is 3-regular.
%\end{theorem}
%
%\begin{proof}
%	From \Cref{isomorfismo.espacio.activo}, we know that $\mathcal{G}(X) \approx \mathcal{G}(X^*)$. Since $\deg(X^*)=3$ and $X^*$ is active, we can apply  \Cref{clasificacion.activo_deg=3} and easily see that $\tau(X^*) \approx X^*$ for every 2-switch $\tau$ and for each of the 11 graphs to which $X^*$ can be isomorphic. Therefore, $\mathcal{G}(X)$ is 3-regular.
%\end{proof}
%
%Regarding  \Cref{2switch.casi-preserva.deg3}, it is important to note that it ceases to be true if $T_{221}$ or $\overline{T_{221}}$ are induced subgraphs of $X$. In fact, in such a case, $X^*\approx T_{221}$ or $\overline{T_{221}}$, and $T_{221}$ has a neighbor of degree 4 ($\overline{R_{211}}$) in $\mathcal{G}(T_{221})$. 
%%(see \Cref{G(T221),split.primos.deg=4.fig.3}). 
%%\begin{figure}[H]
%%	\centering
%%	\begin{tikzpicture}[scale=0.9, every node/.style={scale=0.75}]
%%		% ===== Grafo de realización: columna izquierda =====
%%		\begin{scope}[shift={(-4.5,-0.5)}]
%%			\node[draw, circle, fill=white, thick] (G_1) at (-1,1) {};
%%			\node[color=black,above right] at (G_1) {$G_1$};
%%			\node[draw, circle, fill=white, thick] (G_2) at (1,1) {};
%%			\node[above right] at (G_2) {$G_2$};
%%			\node[draw, circle, fill=white, thick] (G_3) at (-1,-1) {};
%%			\node[below right] at (G_3) {$G_3$};
%%			\node[draw, circle, fill=white, thick] (G_4) at (1,-1) {};
%%			\node[below right] at (G_4) {$G_4$};
%%			\node[color=black,draw, circle, fill=black, thick] (H_1) at (0,0) {};
%%			\node at ($(H_1)+(0,0.4)$) {\textcolor{black}{$H_1$}};
%%			\draw (G_1) -- (G_2);
%%			\draw (G_1) -- (G_3);
%%			\draw (G_2) -- (G_4);
%%			\draw (G_3) -- (G_4);
%%			\draw[thick] (G_1) -- (H_1);
%%			\draw[thick] (G_2) -- (H_1);
%%			\draw[thick] (G_3) -- (H_1);
%%			\draw[thick] (G_4) -- (H_1);
%%		\end{scope}
%%		
%%		% ===== G1: columna centro =====
%%		\begin{scope}[shift={(0,0)}]
%%			% Clique (triángulo pequeño)
%%			\node[draw, circle, fill=white, thick] (G0) at (0,0.9) {};
%%			\node[draw, circle, fill=white, thick] (G1) at (-0.8,-0.45) {};
%%			\node[draw, circle, fill=white, thick] (G2) at (0.8,-0.45) {};
%%			% Independientes rotados 180°
%%			\node[draw, circle, fill=white, thick] (G3) at (0,-1.8) {};
%%			\node[above right] at (G3) {$T_{221}$};
%%			\node[draw, circle, fill=white, thick] (G4) at (1.6,0.9) {};
%%			\node[draw, circle, fill=white, thick] (G5) at (-1.6,0.9) {};
%%			% Clique edges
%%			\draw (G0) -- (G1);
%%			\draw (G1) -- (G2);
%%			\draw (G0) -- (G2);
%%			% (Agregar manualmente las otras aristas)
%%			\draw (G3) -- (G1);
%%			\draw (G4) -- (G2);
%%			\draw (G4) -- (G0);
%%			\draw (G5) -- (G1);
%%			\draw (G5) -- (G0);
%%		\end{scope}
%%		
%%		% ===== H1: columna derecha =====
%%		\begin{scope}[shift={(4.5,0)}]
%%			% Clique (triángulo pequeño)
%%			\node[draw, circle, fill=black, thick] (H0) at (0,0.9) {};
%%			\node[draw, circle, fill=black, thick] (H1) at (-0.8,-0.45) {};
%%			\node[draw, circle, fill=black, thick] (H2) at (0.8,-0.45) {};
%%			% Independientes 
%%			\node[draw, circle, fill=black, thick] (H3) at (0,-1.8) {};
%%			\node[above right] at (H3) {$\overline{R_{211}}$};
%%			\node[draw, circle, fill=black, thick] (H4) at (1.6,0.9) {};
%%			\node[draw, circle, fill=black, thick] (H5) at (-1.6,0.9) {};
%%			% Clique edges
%%			\draw[thick] (H0) -- (H1);
%%			\draw[thick] (H1) -- (H2);
%%			\draw[thick] (H0) -- (H2);
%%			% otras aristas
%%			\draw[thick] (H3) -- (H1);
%%			\draw[thick] (H4) -- (H2);
%%			\draw[thick] (H4) -- (H0);
%%			\draw[thick] (H5) -- (H2);
%%			\draw[thick] (H5) -- (H0);
%%		\end{scope}
%%	\end{tikzpicture}
%%	\caption{Left: $\mathcal{G}(T_{221})$; $G_i\approx T_{221}$ (center), \textcolor{black}{$H_1\approx\overline{R_{211}}$} (right).}
%%	\label{G(T221)}
%%\end{figure}
%
%For every $k\geq 3$ there exist a prime graph $G_k$ with $\deg(G_k)=k$
%and a 2-switch $\tau$ such that $\deg(\tau(G_k))\neq k$. Indeed, let
%$G_k$ be the split graph with clique $K=\{x,y_1,\dots,y_{k-1}\}$,
%independent set $I=\{a,b,c\}$, and $N_a=\{x\}$,
%$N_b=\{y_1,\dots,y_{k-1}\}$, $N_c=\{y_1\}$. By
%\Cref{prop.basicas.sigma_uv}, $\sigma_{ab}=k-1$, $\sigma_{ac}=1$ and
%$\sigma_{bc}=0$, so $\deg(G_k)=k$, and $\Phi(G_k)=cab$ ignoring multiplicities. Every vertex of $K\cup\{a,b\}$ lies in an induced $P_4$ of the form
%$axy_ib$, and $c$ lies in the induced $P_4$ $axy_1c$; hence
%$G_k$ is active, and therefore prime by
%\Cref{S.primo.iff.Phi(S).conexo}. Let
%$\tau$ be the 2-switch replacing $ax,by_1$ with $ay_1,bx$. In
%$\tau(G_k)$ we have $N_a=N_c=\{y_1\}$ and
%$N_b=\{x,y_2,\dots,y_{k-1}\}$, so $\sigma_{ab}=k-1=\sigma_{cb}$ and
%$\sigma_{ac}=0$. Therefore $\deg(\tau(G_k))=2k-2\neq k$. Note that $G_3\approx \overline{T_{221}}$ and $G_4\approx R_{311}$ (see \Cref{grafos.primos.deg=3,split.primos.deg=4.fig.2}). The case
%$k=5$ is shown in \Cref{R.k-1.1.1}.

%
%%For every $k \geq 4$, there exists a graph $G_k$ and a 2-switch $\tau$
%%such that $\deg(\tau(G_k)) \neq \deg(G_k) = k$; that is, the 2-switch
%%does not generally preserve degree $k$ for $k \geq 4$. To see this,
%%suppose that $\tau$ replaces the edges $ab, cd$ with $ad, bc$. If
%%$G_4 = abxcd \approx P_5$, then $\tau(G_4) \approx K_3 \,\dot\cup\, K_2$
%%and therefore $\deg(\tau(G_4)) = 6 \neq 4 = \deg(G_4)$. For $k \geq 5$,
%%consider $G_k = \overline{D_{1,k-4}} \circ G_4$. We know that
%%$\deg(\overline{D_{1,k-4}}) = k-4$. Hence, by
%%\Cref{hechos.previos},
%%\[
%%\deg(G_k) = k
%%\quad\text{and}\quad
%%\deg(\tau(G_k)) = \deg(\overline{D_{1,k-4}} \circ \tau(G_4)) = k+2.
%%\] 
%
%The classification obtained in the previous sections determines the
%realization graph of every prime graph of degree at most $4$. Since $\mathcal{G}(d)\approx\mathcal{G}(\overline{d})$ by
%\Cref{dual.spaces.iso}, only one sequence in each complementary pair
%has to be treated. 
%
%\begin{theorem}
%	\label{realization.graphs.deg<5}
%	Let $X$ be a prime graph with $\deg(X)\leq 4$. Then $\mathcal{G}(X)$
%	is isomorphic to one of the ten graphs of
%	\Cref{realization.graphs.fig}, according to the following table.
%	\begin{center}
%		\begin{tabular}{ll}
%			\hline
%			$\mathcal{G}(X)$ & $X$ \\
%			\hline
%			$K_2$ & $P_4$\\
%			$K_3$ & $C_4$, $2K_2$, $D_{2,1}$, $\overline{D_{2,1}}$\\
%			$K_4$ & $D_{3,1}$, $\overline{D_{3,1}}$\\
%			$K_{3,3}$ & $T_{111}$, $\overline{T_{111}}$\\
%			$K_{1,2,2}$ & $T_{221}$, $\overline{T_{221}}$, $R_{211}$, $\overline{R_{211}}$\\
%			$K_5$ & $D_{4,1}$, $\overline{D_{4,1}}$\\
%			$K_{2,2,2}$ & $D_{2,2}$, $\overline{D_{2,2}}$, $B_1$, $\overline{B_1}$, $B_3$, $\overline{B_3}$\\
%			$K_{1,3,3}$ & $B_2$, $\overline{B_2}$\\
%			$\mathcal{G}_1$ & $R_{311}$, $\overline{R_{311}}$, $R_{331}$, $\overline{R_{331}}$\\
%			$\mathcal{G}_2$ & $R_{321}$, $\overline{R_{321}}$\\
%			\hline
%		\end{tabular}
%	\end{center}
%\end{theorem}
%
%\begin{proof}
%	By
%	\Cref{clasificacion.activo_deg=2,clasificacion.activo_deg=3,clasificacion.split.primos.deg=4,clasificacion.no.split.primos.deg=4}
%	there are exactly $29$ prime graphs of degree at most $4$, and they
%	realize $26$ distinct degree sequences: the only coincidences are
%	$d(T_{221})=d(\overline{R_{211}})$, $d(\overline{T_{221}})=d(R_{211})$
%	and $d(B_1)=d(\overline{B_1})$. For each of these sequences, listing
%	the realizations of $d$ and the 2-switches between them is a finite
%	computation, which yields the table.
%\end{proof}
%\begin{figure}[H]
%	\centering
%	\begin{tikzpicture}[scale=0.85,
%		every node/.style={draw, circle, fill=white, inner sep=0pt, minimum size=6.5pt},
%		nm/.style={draw=none, fill=none},
%		]
%		
%		%%%%%%%%%%%%%%%%%%%% row 1 %%%%%%%%%%%%%%%%%%%%
%		\begin{scope}[shift={(0,0)}]
%			\node (a) at (0.15,0.9) {};
%			\node (b) at (1.65,0.9) {};
%			\draw (a)--(b);
%			\node[nm] at (0.9,-0.45) {$K_2$};
%		\end{scope}
%		
%		\begin{scope}[shift={(3.4,0)}]
%			\node (a) at (0.9,1.7) {};
%			\node (b) at (0.15,0.45) {};
%			\node (c) at (1.65,0.45) {};
%			\draw (a)--(b)--(c)--(a);
%			\node[nm] at (0.9,-0.45) {$K_3$};
%		\end{scope}
%		
%		\begin{scope}[shift={(6.8,0)}]
%			\node (a) at (0.2,1.65) {};
%			\node (b) at (1.6,1.65) {};
%			\node (c) at (1.6,0.25) {};
%			\node (d) at (0.2,0.25) {};
%			\draw (a)--(b)--(c)--(d)--(a); \draw (a)--(c); \draw (b)--(d);
%			\node[nm] at (0.9,-0.45) {$K_4$};
%		\end{scope}
%		
%		\begin{scope}[shift={(10.2,0)}]
%			\node (a) at (0.90,1.78) {};
%			\node (b) at (0.09,1.19) {};
%			\node (c) at (0.40,0.24) {};
%			\node (d) at (1.40,0.24) {};
%			\node (e) at (1.71,1.19) {};
%			\draw (a)--(b)--(c)--(d)--(e)--(a);
%			\draw (a)--(c); \draw (a)--(d); \draw (b)--(d); \draw (b)--(e); \draw (c)--(e);
%			\node[nm] at (0.9,-0.45) {$K_5$};
%		\end{scope}
%		
%		%%%%%%%%%%%%%%%%%%%% row 2 %%%%%%%%%%%%%%%%%%%%
%		\begin{scope}[shift={(0,-3.3)}]
%			\node (t1) at (0.15,1.65) {};
%			\node (t2) at (0.90,1.65) {};
%			\node (t3) at (1.65,1.65) {};
%			\node (b1) at (0.15,0.25) {};
%			\node (b2) at (0.90,0.25) {};
%			\node (b3) at (1.65,0.25) {};
%			\draw (t1)--(b1); \draw (t1)--(b2); \draw (t1)--(b3);
%			\draw (t2)--(b1); \draw (t2)--(b2); \draw (t2)--(b3);
%			\draw (t3)--(b1); \draw (t3)--(b2); \draw (t3)--(b3);
%			\node[nm] at (0.9,-0.45) {$K_{3,3}$};
%		\end{scope}
%		
%		\begin{scope}[shift={(3.4,-3.3)}]
%			\node (a) at (0.15,1.65) {};
%			\node (b) at (1.65,1.65) {};
%			\node (c) at (1.65,0.25) {};
%			\node (d) at (0.15,0.25) {};
%			\node (o) at (0.90,0.95) {};
%			\draw (a)--(b)--(c)--(d)--(a);
%			\draw (o)--(a); \draw (o)--(b); \draw (o)--(c); \draw (o)--(d);
%			\node[nm] at (0.9,-0.45) {$K_{1,2,2}$};
%		\end{scope}
%		
%		\begin{scope}[shift={(6.8,-3.3)}]
%			\node (h1) at (0.90,1.80) {};
%			\node (h2) at (0.12,1.35) {};
%			\node (h3) at (0.12,0.45) {};
%			\node (h4) at (0.90,0.00) {};
%			\node (h5) at (1.68,0.45) {};
%			\node (h6) at (1.68,1.35) {};
%			\draw (h1)--(h2)--(h3)--(h4)--(h5)--(h6)--(h1);
%			\draw (h1)--(h3); \draw (h2)--(h4); \draw (h3)--(h5);
%			\draw (h4)--(h6); \draw (h5)--(h1); \draw (h6)--(h2);
%			\node[nm] at (0.9,-0.45) {$K_{2,2,2}$};
%		\end{scope}
%		
%		\begin{scope}[shift={(10.2,-3.3)}]
%			\node (h1) at (0.90,1.80) {};
%			\node (h2) at (0.12,1.35) {};
%			\node (h3) at (0.12,0.45) {};
%			\node (h4) at (0.90,0.00) {};
%			\node (h5) at (1.68,0.45) {};
%			\node (h6) at (1.68,1.35) {};
%			\node (o)  at (0.90,0.90) {};
%			\draw (h1)--(h2)--(h3)--(h4)--(h5)--(h6)--(h1);
%			\draw (o)--(h1); \draw (o)--(h2); \draw (o)--(h3);
%			\draw (o)--(h4); \draw (o)--(h5); \draw (o)--(h6);
%			\draw (h1) to[bend left=55] (h4);
%			\draw (h2) to[bend right=55] (h5);
%			\draw (h3) to[bend left=55] (h6);
%			\node[nm] at (0.9,-0.45) {$K_{1,3,3}$};
%		\end{scope}
%		
%		%%%%%%%%%%%%%%%%%%%% row 3 %%%%%%%%%%%%%%%%%%%%
%		\begin{scope}[shift={(3.4,-6.6)}]
%			\node (h1) at (0.90,1.80) {};
%			\node (h2) at (0.12,1.35) {};
%			\node (h3) at (0.12,0.45) {};
%			\node (h4) at (0.90,0.00) {};
%			\node (h5) at (1.68,0.45) {};
%			\node (h6) at (1.68,1.35) {};
%			\node (o)  at (0.90,0.90) {};
%			\draw (h1)--(h2)--(h3)--(h4)--(h5)--(h6)--(h1);
%			\draw (o)--(h1); \draw (o)--(h2); \draw (o)--(h3);
%			\draw (o)--(h4); \draw (o)--(h5); \draw (o)--(h6);
%			\draw (h1)--(h3); \draw (h4)--(h6);
%			\draw (h2) to[bend right=55] (h5);
%			\node[nm] at (0.9,-0.45) {$\mathcal{G}_1$};
%		\end{scope}
%		
%		\begin{scope}[shift={(7.4,-6.6)}]
%			\node (c1) at (0.00,1.90) {};
%			\node (c2) at (1.90,1.90) {};
%			\node (c3) at (1.90,0.00) {};
%			\node (c4) at (0.00,0.00) {};
%			\node (k1) at (0.95,1.45) {};
%			\node (k2) at (1.60,0.95) {};
%			\node (k3) at (0.95,0.45) {};
%			\node (k4) at (0.30,0.95) {};
%			\draw (c1)--(c2)--(c3)--(c4)--(c1);
%			\draw (k1)--(k2)--(k3)--(k4)--(k1); \draw (k1)--(k3); \draw (k2)--(k4);
%			\draw (c1)--(k1); \draw (c2)--(k1); \draw (c2)--(k2); \draw (c3)--(k2);
%			\draw (c3)--(k3); \draw (c4)--(k3); \draw (c4)--(k4); \draw (c1)--(k4);
%			\node[nm] at (0.95,-0.45) {$\mathcal{G}_2$};
%		\end{scope}
%		
%	\end{tikzpicture}
%	\caption{The ten realization graphs of \Cref{realization.graphs.deg<5}.}
%	\label{realization.graphs.fig}
%\end{figure}

%%%%%%%%%%%%%%%%%%%%%%%%%%%%%%%%%%%%%%%%%%%%%%%%%%%%%%%%%%%%%%%%%%%%%%%%
%%%%%%%%%%%%%%%%%%%%%%%%%%%%%%%%%%%%%%%%%%%%%%%%%%%%%%%%%%%%%%%%%%%%

\section*{Acknowledgements}
%The authors would like to thank the anonymous referees, whose comments greatly improved the presentation of this manuscript.
This work was partially supported by Universidad Nacional de San Luis, grants PI UNSL 03-1326, PROIPRO 03-2923, and PROINI 03-124, and Consejo Nacional de Investigaciones Cient\'ificas y T\'ecnicas grant PIP $11220220$ $100068$CO. 

%%%%%%%%%%%%%%%%%%%%%%%%%%%%%%%%%%%%%%%%%%%%%%%%%%%%%%%%%%%%%%%%


\begin{thebibliography}{99}
	
	\bibitem{arikati1999realization}
	S.~R. Arikati and U.~N. Peled.
	\newblock The realization graph of a degree sequence with majorization gap 1 is
	hamiltonian.
	\newblock {\em Linear Algebra and its Applications}, 290(1-3):213--235, 1999.
	
	\bibitem{barrus2016realization}
	M.~D. Barrus.
	\newblock On realization graphs of degree sequences.
	\newblock {\em Discrete Mathematics}, 339(8):2146--2152, 2016.
	
	\bibitem{barrus2023cliques}
	M.~D. Barrus and N.~Haronian.
	\newblock Cliques in realization graphs.
	\newblock {\em Discrete Mathematics}, 346(1):113184, 2023.
	
	\bibitem{barrus.west.A4}
	M.~D. Barrus and D.~B. West.
	\newblock The {$A_4$}-structure of a graph.
	\newblock {\em Journal of Graph Theory}, 69(2):97--113, 2012.
	
	\bibitem{barrus2012switches}
	M.~D. Barrus.
	\newblock On 2-switches and isomorphism classes.
	\newblock {\em Discrete Mathematics}, 312(15):2217--2222, 2012.
	
%	\bibitem{barrus2016realization}
%	M.~D. Barrus.
%	\newblock On realization graphs of degree sequences.
%	\newblock {\em Discrete Mathematics}, 339(8):2146--2152, 2016.
	
	
	\bibitem{chartrand2010graphs}
	G.~Chartrand, L.~Lesniak, and P.~Zhang.
	\newblock {\em Graphs \& Digraphs}.
	\newblock Textbooks in Mathematics. CRC Press, 5th edition, 2010.
	
	\bibitem{splitnordhausgaddum}
	C.~Cheng, K.~L. Collins, and A.~N. Trenk.
	\newblock Split graphs and {N}ordhaus--{G}addum graphs.
	\newblock {\em Discrete Mathematics}, 339(9):2345--2356, 2016.
	
%	\bibitem{eggleton2006realizations}
%	R.~B. Eggleton and D.~A. Holton.
%	\newblock Simple and multigraphic realizations of degree sequences.
%	\newblock In {\em Combinatorial Mathematics VIII}, volume 884 of {\em Lecture
%		Notes in Mathematics}, pages 155--172. Springer, 1981.
	
	\bibitem{hammer1977split}
	S.~F{\"o}ldes and P.~L. Hammer.
	\newblock Split graphs.
	\newblock {\em Congressus Numerantium}, 19:311--315, 1977.
	
%	\bibitem{fulkerson1965}
%	D.~R. Fulkerson, A.~J. Hoffman, and M.~H. McAndrew.
%	\newblock Some properties of graphs with multiple edges.
%	\newblock {\em Canadian Journal of Mathematics}, 17:166--177, 1965.
	
	\bibitem{gallai1967}
	T.~Gallai.
	\newblock Transitiv orientierbare {G}raphen.
	\newblock {\em Acta Mathematica Academiae Scientiarum Hungaricae},
	18:25--66, 1967.
	
		\bibitem{jaume2025nullspace}
	D.~A. Jaume, V.~N. Schv{\"o}llner, C.~Panelo, and K.~Pereyra.
	\newblock On the nullspace of split graphs.
	\newblock {\em Linear Algebra and its Applications}, 747:98--125, 2026.
	
	
%	\bibitem{mahadev1995threshold}
%	N.~V.~R. Mahadev and U.~N. Peled.
%	\newblock {\em Threshold Graphs and Related Topics}, volume~56 of {\em Annals of
%		Discrete Mathematics}.
%	\newblock Elsevier, 1995.
	
	\bibitem{pastine2025simple}
	V.~N. Schv{\"o}llner and A.~Pastine.
	\newblock Simple factor graphs associated with split graphs.
	\newblock {\em arXiv preprint arXiv:2512.24252}, 2025.
	
	
	\bibitem{pastine20252}
	V.~N. Schv{\"o}llner and A.~Pastine.
	\newblock The 2-switch-degree of a graph.
	\newblock {\em arXiv preprint arXiv:2511.23327}, 2025.
	
	
	\bibitem{schvollner2026induced}
	V.~N. Schv{\"o}llner and A.~Pastine.
	\newblock Induced paths and cycles in factor graphs of split graphs.
	\newblock {\em arXiv preprint arXiv:2603.14061}, 2026.
	
	\bibitem{schvollner_pseudoforests}
	V.~N. Schv{\"o}llner, A.~Pastine, and D.~A. Jaume.
	\newblock 2-switch: transition and stability on forests and pseudoforests.
		\newblock {\em Bulletin of the Institute of Combinatorics and its
		Applications}, 110, 2027. In press. Preprint arXiv:2603.07439.
	
	\bibitem{taylor2006contrained}
	R.~Taylor.
	\newblock Constrained switchings in graphs.
	\newblock In {\em Combinatorial Mathematics VIII}, volume 884 of {\em Lecture
		Notes in Mathematics}, pages 314--336. Springer, 1981.
	
	\bibitem{tyshkevich2000decomposition}
	R.~Tyshkevich.
	\newblock Decomposition of graphical sequences and unigraphs.
	\newblock {\em Discrete Mathematics}, 220(1-3):201--238, 2000.
	
	\bibitem{whitman2020split}
	R.~Whitman.
	\newblock Split graphs, unigraphs, and {T}yshkevich decompositions.
	\newblock Undergraduate honors thesis, Wellesley College, 2020.
	
\end{thebibliography}
	
%\bibliographystyle{abbrv}
%\bibliography{citas__classification_deg_less_than_5}	
	

\end{document}
